% AVISO: este archivo incluye la sección «Análisis exegético pasaje por pasaje» (11:1-36, «ahora» de 11:30-31, «y así», «de Sion») y la predicación «El misterio revelado». Editar siempre sobre main actualizado (git pull); no regenerar desde copias antiguas.
\title[Israel 8: Romanos 11]{\bfseries{Israel, el pueblo de Dios}}
\subtitle{Estudio 8. Romanos 11: la restauración de Israel}
\author{}\institute{}\date{}
\titlegraphic{\includegraphics[width=0.85\linewidth,height=0.38\textheight,keepaspectratio]{portada.jpg}\\[-1pt]
  {\fontsize{5}{6}\selectfont\color{sepia}\textit{Nota.} Elaboración propia: Ro 11:26a en escritura uncial (texto NA28).}}
\AtBeginSection[]{\begin{frame}\frametitle{Contenidos}\tableofcontents[currentsection]\end{frame}}
\newcommand{\mpt}[2]{({(#1-12)*0.26},{(#2-30)*0.33})}
\begin{document}
\frame{\titlepage}
\begin{frame}\frametitle{Contenidos}\tableofcontents\end{frame}

% ===================================================
\section{Introducción}
\begin{frame}{¿Ha desechado Dios a su pueblo?}
\begin{columns}[T,onlytextwidth]
\begin{column}{0.52\textwidth}
\begin{cita}Digo, pues: ¿Ha desechado Dios a su pueblo? ¡De ningún modo! Porque también yo soy israelita, de la descendencia de Abraham, de la tribu de Benjamín. No ha desechado Dios a su pueblo, al cual conoció de antemano.\fuente{(Ro 11:1--2)}\end{cita}
\begin{cita}Porque irrevocables son los dones y el llamamiento de Dios.\fuente{(Ro 11:29)}\end{cita}
\end{column}
\begin{column}{0.45\textwidth}\small
Roma, c. 57 e.\,c.: la mayoría de Israel no cree en Jesús; crecen las iglesias de gentiles. Conclusión rápida: \emph{Dios cambió de pueblo}. Pablo la desmonta.\par\medskip
\destacar{Tesis 1. Identidad y vocación.} La vocación de Israel se preserva en un remanente y alcanzará su plenitud (11:5, 12, 29).\par\medskip
\destacar{Tesis 2. Restauración.} El endurecimiento es parcial y temporal; «y así todo Israel será salvo» (11:25--27).
\end{column}
\end{columns}
\end{frame}

\begin{frame}{La culminación de la serie}
\centering
\begin{tikzpicture}
 \foreach \x/\y/\n/\t/\r in {0/0/1/{Abraham y el pacto}/{Raíz santa; «amados por causa de los padres» (11:16, 28)},
   3.15/0/2/{Sinaí y vocación}/{Celos por «lo que no es pueblo» (Dt 32:21; 11:11)},
   6.3/0/3/{Jueces y Samuel}/{«YHWH no desamparará a su pueblo» (1~S 12:22; 11:1)},
   9.45/0/4/{David y la división}/{David testigo (11:9); raíz de Isaí (15:12)},
   1.6/-2.7/5/{Reyes y cautiverios}/{Los 7000 de Elías; nuevo pacto (11:4, 27)},
   4.75/-2.7/6/{Segundo Templo}/{«Vendrá de Sion el Redentor» (Is 59:20; 11:26)},
   7.9/-2.7/7/{Cristo y el año 70}/{Remanente; «hasta que» (Lc 21:24; 11:25)}}{
  \node[caja,anchor=north] (b\n) at (\x,\y) {\parbox[t][1.95cm][t]{2.65cm}{\RaggedRight
    {\small\textbf{\textcolor{ladrillo}{\t}}}\\[3pt]\scriptsize \r}};
  \node[circle,fill=cafe,text=piedra,font=\scriptsize\bfseries,inner sep=1.6pt] at (b\n.north west) {\n};}
\end{tikzpicture}\par\vspace{4pt}
{\small\itshape\color{sepia}Romanos 11 no inventa una teología de Israel: recoge la historia de los siete estudios anteriores.}
\end{frame}

\begin{frame}{Cuatro dimensiones}
\centering
\begin{tikzpicture}
 \foreach \x/\n/\t/\d/\r in {0/1/{Histórica}/{¿Qué situación de Roma y de Israel explica Ro 11?}/{Hch 18:2; 21:20; Ro 15:25--27; Suetonio},
   3.15/2/{Literaria}/{¿Cómo está construido el argumento de Ro 9--11?}/{Diatriba; citas del AT; olivo; doxología},
   6.3/3/{Filológica}/{¿Qué dicen exactamente las palabras griegas?}/{NA28; BHS; LXX (Rahlfs)},
   9.45/4/{Teológica}/{¿Qué afirma sobre elección, remanente y salvación de Israel?}/{Ro 11:1--2, 5, 11--15, 25--32}}{
  \node[caja,anchor=north] (b\n) at (\x,0) {\parbox[t][2.9cm][t]{2.65cm}{\RaggedRight
    {\small\textbf{\textcolor{ladrillo}{\t}}}\\[4pt]\small \d\par\vfill{\tiny\color{sepia}(\r)}}};
  \node[circle,fill=cafe,text=piedra,font=\scriptsize\bfseries,inner sep=1.6pt] at (b\n.north west) {\n};}
\end{tikzpicture}\par\vspace{6pt}
{\scriptsize\color{sepia}Las citas bíblicas son traducción propia desde el texto hebreo (BHS) y griego (NA28); LXX según Rahlfs-Hanhart. Numeración española [hebrea o griega] cuando difiere. En Ro 9--11, «Israel» es siempre el pueblo completo, no el antiguo reino del norte.}
\end{frame}

% ===================================================
\section{Historia}
\begin{frame}{Roma y la carta: cronología}
\centering
\begin{tikzpicture}[x=0.82cm]
\draw[line width=1.6pt,cafe,-{Latex[length=3mm]}] (0,0) -- (15.6,0);
\foreach \x/\t/\b in {0.5/{Pompeyo toma Jerusalén}/{63 a.\,e.\,c.; cautivos a Roma},
                      2.2/{Expulsión bajo Tiberio}/{19 e.\,c.; Josefo, \emph{Ant.} 18.81--84},
                      3.9/{Edicto de Claudio}/{c. 49; Suetonio; Hch 18:2},
                      5.6/{Pablo en Corinto}/{c. 50--52; Hch 18},
                      7.3/{Muere Claudio}/{54; regreso (Ro 16:3)},
                      9.0/{Carta a los Romanos}/{c. 56--57, desde Corinto},
                      10.7/{Colecta y arresto}/{c. 57; Hch 21},
                      12.4/{Persecución de Nerón}/{64; Tácito, \emph{Ann.} 15.44},
                      14.1/{Caída del templo}/{70 (Estudio 7)}}{
  \fill[oro] (\x,0) circle (4pt);
  \node[above,font=\scriptsize\bfseries,text=ladrillo,text width=1.25cm,align=center] at (\x,0.2) {\t};
  \node[below,font=\tiny,text=sepia,text width=1.25cm,align=center] at (\x,-0.2) {\b};}
\draw[decorate,decoration={brace,amplitude=5pt,mirror},sepia] (3.6,-1.2) -- (7.6,-1.2) node[midway,below=5pt,font=\scriptsize\itshape]{expulsión y regreso};
\draw[decorate,decoration={brace,amplitude=5pt,mirror},sepia] (8.7,-1.2) -- (11.0,-1.2) node[midway,below=5pt,font=\scriptsize\itshape]{Ro 15:25--32};
\end{tikzpicture}\par\vspace{6pt}
\begin{minipage}{0.95\linewidth}\precision{Galión, procónsul de Acaya c. 51--52, está atestiguado por una inscripción de Delfos. La identificación del Erasto de Ro 16:23 con el edil de la inscripción de Corinto está en debate.}\end{minipage}\par
{\fontsize{5}{6}\selectfont\color{sepia}\raggedleft\textit{Nota.} Elaboración propia. Fechas aproximadas.\par}
\end{frame}

\begin{frame}{Corinto, Roma y Jerusalén}
\begin{columns}[c,onlytextwidth]
\begin{column}{0.66\textwidth}\centering\includegraphics[width=\linewidth,height=0.86\textheight,keepaspectratio]{mapa_romanos.pdf}\par\vspace{1pt}{\fontsize{4.5}{5.5}\selectfont\color{sepia}\textit{Nota.} Elaboración propia (GSHHS/WDB, Basemap); ubicaciones aproximadas.\par}\end{column}
\begin{column}{0.32\textwidth}\scriptsize\RaggedRight
{\scriptsize\RaggedRight
\definecolor{mapacarmesi}{HTML}{D62839}
\definecolor{mapanaranja}{HTML}{F28C28}
\definecolor{mapavioleta}{HTML}{7B2CBF}
\definecolor{mapaturquesa}{HTML}{0F8B8D}
\definecolor{mapaazul}{HTML}{1D4E89}
\definecolor{mapaocre}{HTML}{B5651D}
\tikz[baseline=-0.6ex]{\draw[white,line width=3.2pt,line cap=round](0,0)--(0.62,0);\draw[mapavioleta,line width=1.8pt,line cap=round,dash pattern=on 1.4pt off 1.6pt,-{Latex[length=1.6mm]}](0,0)--(0.62,0);}~Edicto de Claudio: Aquila y Priscila (Hch 18:2)\par
\tikz[baseline=-0.6ex]{\draw[white,line width=3.2pt,line cap=round](0,0)--(0.62,0);\draw[mapacarmesi,line width=1.8pt,line cap=round,-{Latex[length=1.6mm]}](0,0)--(0.62,0);}~La carta con Febe: Cencrea $\to$ Roma (Ro 16:1)\par
\tikz[baseline=-0.6ex]{\draw[white,line width=3.2pt,line cap=round](0,0)--(0.62,0);\draw[mapanaranja,line width=1.8pt,line cap=round,dash pattern=on 3.6pt off 2pt,-{Latex[length=1.6mm]}](0,0)--(0.62,0);}~La colecta: Corinto $\to$ Jerusalén (Ro 15:25–26; Hch 20–21)\par
\tikz[baseline=-0.6ex]{\draw[white,line width=3.2pt,line cap=round](0,0)--(0.62,0);\draw[mapacarmesi,line width=1.8pt,line cap=round,dash pattern=on 1.4pt off 1.6pt,-{Latex[length=1.6mm]}](0,0)--(0.62,0);}~Proyecto: hacia España (Ro 15:24, 28)\par
\smallskip{\color{black!70}\textbf{Itinerario}}\par
\tikz[baseline=-0.7ex]\node[circle,fill=mapacarmesi,draw=white,line width=0.5pt,text=white,font=\fontsize{5}{5}\selectfont\bfseries,inner sep=0.6pt,minimum size=2.6mm]{1};\,Cencrea\quad \tikz[baseline=-0.7ex]\node[circle,fill=mapacarmesi,draw=white,line width=0.5pt,text=white,font=\fontsize{5}{5}\selectfont\bfseries,inner sep=0.6pt,minimum size=2.6mm]{2};\,Roma\quad \tikz[baseline=-0.7ex]\node[circle,fill=mapanaranja,draw=white,line width=0.5pt,text=white,font=\fontsize{5}{5}\selectfont\bfseries,inner sep=0.6pt,minimum size=2.6mm]{3};\,Filipos\quad \tikz[baseline=-0.7ex]\node[circle,fill=mapanaranja,draw=white,line width=0.5pt,text=white,font=\fontsize{5}{5}\selectfont\bfseries,inner sep=0.6pt,minimum size=2.6mm]{4};\,Mileto\quad \tikz[baseline=-0.7ex]\node[circle,fill=mapanaranja,draw=white,line width=0.5pt,text=white,font=\fontsize{5}{5}\selectfont\bfseries,inner sep=0.6pt,minimum size=2.6mm]{5};\,Jerusalén\par}

\end{column}
\end{columns}
\end{frame}

\begin{frame}{Corinto, Roma y Jerusalén: claves}
\small\RaggedRight
Pablo escribe desde \textbf{Corinto} (Hch 20:2--3; Ro 16:1, 23) a una iglesia que \textbf{no fundó}.\par\medskip
La carta viaja con Febe; él lleva la colecta a Jerusalén y después piensa ir a Roma y a España (Ro 15:24--28).\par\medskip
Aquila y Priscila llegaron a Corinto por el edicto de Claudio (Hch 18:2).
\end{frame}

\begin{frame}{El edicto de Claudio}
\begin{columns}[T,onlytextwidth]
\begin{column}{0.50\textwidth}
\begin{cita}A los judíos, que provocaban continuos tumultos por instigación de Cresto, los expulsó de Roma.\fuente{(Suetonio, \emph{Claudio} 25.4)}\end{cita}
\begin{cita}Encontró a un judío llamado Aquila [\ldots] recién venido de Italia con Priscila su mujer, porque Claudio había mandado que todos los judíos salieran de Roma.\fuente{(Hch 18:2)}\end{cita}
\end{column}
\begin{column}{0.47\textwidth}\small\RaggedRight
\begin{itemize}
\item \emph{Chrestus} y \emph{Christus} sonaban casi igual: probablemente, disputas en las sinagogas sobre Jesús como Mesías.
\item Fecha: c. 49 (Orosio 7.6.15); Dión Casio 60.6.6 habla de una prohibición de reunirse (41).
\item 49--54: iglesias casi solo de gentiles. 54: muere Claudio; los judíos vuelven (Ro 16:3--5).
\end{itemize}
\precision{la hipótesis de Wiefel (1991) explica así la arrogancia gentil que Pablo corrige; es discutida en sus detalles.}
\end{column}
\end{columns}
\end{frame}

\begin{frame}{Remanente, colecta y arrogancia}
\begin{columns}[T,onlytextwidth]
\begin{column}{0.49\textwidth}
\begin{cita}Ya ves, hermano, cuántas \textbf{miríadas} de judíos hay que han creído, y todos son celosos de la ley.\fuente{(Hch 21:20)}\end{cita}
\begin{cita}[fontupper=\scriptsize\itshape]Si los gentiles han sido hechos \textbf{participantes} de sus bienes espirituales, deben también servirles con los materiales.\fuente{(Ro 15:27)}\end{cita}
\small \grc{ἐκοινώνησαν} (15:27) anticipa \grc{συγκοινωνός} del olivo (11:17).
\end{column}
\begin{column}{0.49\textwidth}\small
\textbf{\textcolor{cafe}{«A vosotros hablo, gentiles»}} (11:13)
\begin{itemize}
\item «No te jactes contra las ramas» (11:18).
\item «No te ensoberbezcas, sino teme» (11:20).
\item «Para que no seáis sabios en vuestra propia opinión» (11:25).
\end{itemize}
\textbf{\textcolor{cafe}{Romanos en contexto judío}}: Stendahl (1976), Ro 9--11 es el clímax; Sanders (1977); Nanos (1996); Das (2007).\par\smallskip
\destacar{Ro 11 es corrección pastoral, no curiosidad escatológica.}
\end{column}
\end{columns}
\end{frame}

% ===================================================
\section{Literatura}
\begin{frame}{Mapa de Romanos 9--11}
\centering
\begin{tikzpicture}[node distance=0.25cm]
\node[caja,text width=1.9cm,align=center,minimum height=1.9cm] (a) {\textbf{\textcolor{ladrillo}{Lamento}}\\[2pt]9:1--5\\[2pt]{\tiny privilegios de Israel; «bendito por los siglos, amén»}};
\node[caja,text width=2.1cm,align=center,minimum height=1.9cm,right=of a] (b) {\textbf{\textcolor{ladrillo}{La palabra no ha fallado}}\\[2pt]9:6--29\\[2pt]{\tiny elección; remanente (Is 10:22)}};
\node[caja,text width=2.1cm,align=center,minimum height=1.9cm,right=of b] (c) {\textbf{\textcolor{ladrillo}{La justicia de la fe}}\\[2pt]9:30--10:21\\[2pt]{\tiny tropiezo; «todo el que invoque»; celos}};
\node[caja,text width=2.1cm,align=center,minimum height=1.9cm,right=of c] (d) {\textbf{\textcolor{ladrillo}{No ha desechado}}\\[2pt]11:1--32\\[2pt]{\tiny remanente; olivo; misterio}};
\node[caja,text width=1.9cm,align=center,minimum height=1.9cm,right=of d,fill=oro!25] (e) {\textbf{\textcolor{ladrillo}{Doxología}}\\[2pt]11:33--36\\[2pt]{\tiny «a él la gloria, amén»}};
\foreach \u/\v in {a/b,b/c,c/d,d/e}{\draw[-{Latex[length=2mm]},oro!90!ladrillo,line width=1.2pt] (\u) -- (\v);}
\draw[sepia,line width=0.8pt,-{Latex[length=2mm]}] (a.north) to[bend left=22] node[midway,above,font=\scriptsize\itshape]{inclusión: bendición con «amén» (9:5) — doxología con «amén» (11:36)} (e.north);
\node[below=0.25cm of b,font=\scriptsize,text=cafe,text width=2.3cm,align=center] {¿Por qué no creen? \emph{Elección}};
\node[below=0.25cm of c,font=\scriptsize,text=cafe,text width=2.3cm,align=center] {¿Son culpables? \emph{Responsabilidad}};
\node[below=0.25cm of d,font=\scriptsize,text=cafe,text width=2.3cm,align=center] {¿Qué hará Dios? \emph{Restauración}};
\end{tikzpicture}\par\vspace{6pt}
{\small\itshape\color{sepia}Leer Ro 9:6 («no todos los de Israel son Israel») sin Ro 11:26 («todo Israel será salvo») es leer la mitad del argumento.}\par
{\fontsize{5}{6}\selectfont\color{sepia}\raggedleft\textit{Nota.} Elaboración propia: estructura de Ro 9--11.\par}
\end{frame}

\begin{frame}{Diatriba y argumento}
\begin{columns}[T,onlytextwidth]
\begin{column}{0.42\textwidth}\small\RaggedRight
\textbf{\textcolor{cafe}{\grc{μὴ γένοιτο}: «¡de ningún modo!»}}\par\smallskip
Fórmula de la diatriba (Stowers, 1981); diez veces en Romanos (conteo propio).
\begin{itemize}
\item \grc{μὴ ἀπώσατο ὁ θεὸς τὸν λαὸν αὐτοῦ;} (11:1)
\item \grc{μὴ ἔπταισαν ἵνα πέσωσιν;} (11:11)
\end{itemize}
\grc{μή} anticipa respuesta negativa.\par\smallskip
Cambio de interlocutor: Israel en 3.ª persona (11:1--12) → «a vosotros, gentiles» (11:13) → «tú, olivo silvestre» (11:17).
\end{column}
\begin{column}{0.56\textwidth}
\centering
\begin{tikzpicture}[y=-1.05cm]
\foreach \i/\t/\r/\d in {0/{1. Remanente}/{11:1--10}/{presente: un resto escogido por gracia},
                        1/{2. Tropiezo, no caída}/{11:11--24}/{«cuánto más»; el olivo},
                        2/{3. Misterio}/{11:25--32}/{futuro: todo Israel será salvo},
                        3/{4. Doxología}/{11:33--36}/{el argumento se vuelve alabanza}}{
  \node[caja,text width=5.6cm,minimum height=0.8cm] (s\i) at (0,\i) {\textbf{\textcolor{ladrillo}{\t}} {\tiny\color{sepia}(\r)}\\{\scriptsize \d}};}
\foreach \a/\b in {0/1,1/2,2/3}{\draw[-{Latex[length=2mm]},oro!90!ladrillo,line width=1.2pt] (s\a) -- (s\b);}
\end{tikzpicture}
\end{column}
\end{columns}
\end{frame}

\begin{frame}{Las citas del AT en Ro 11}
\centering\tiny\renewcommand{\arraystretch}{1.2}
\begin{tabular}{>{\RaggedRight\arraybackslash}p{1.9cm}>{\RaggedRight\arraybackslash}p{2.9cm}>{\RaggedRight\arraybackslash}p{2.7cm}>{\RaggedRight\arraybackslash}p{3.6cm}}
\toprule
\textbf{Ro 11 / fuente} & \textbf{TM} & \textbf{LXX} & \textbf{Pablo}\\\midrule
11:4; 1~R 19:18 & «Dejaré (\heb{וְהִשְׁאַרְתִּי}) siete mil» & «Dejarás siete mil» & \grc{κατέλιπον ἐμαυτῷ}: más cerca del TM; «para mí»\\
11:8; Dt 29:4 [29:3] + Is 29:10 & «no os dio corazón»; \heb{רוּחַ תַּרְדֵּמָה} & \grc{πνεῦμα κατανύξεως} & Funde dos textos: «les dio ojos para no ver»\\
11:9--10; Sal 69:22--23 [69:23--24; LXX 68:23--24] & «Sea su mesa un lazo» & \grc{εἰς παγίδα\ldots} & Sigue la LXX; añade \grc{εἰς θήραν}\\
11:26--27; Is 59:20--21 + 27:9 & «vendrá \emph{a} Sion» & «\emph{por causa de} Sion» & «\emph{de} Sion»; Is 27:9 LXX; eco de Jer 31:33--34\\
11:34; Is 40:13 & «¿quién midió el espíritu de YHWH?» & \grc{τίς ἔγνω νοῦν κυρίου;} & LXX casi literal\\
11:35; Job 41:11 [41:3] & «¿quién se me adelantó?» & otra lectura & Sigue el sentido del TM\\
\bottomrule
\end{tabular}\par\vspace{6pt}
\begin{minipage}{0.95\linewidth}\precision{\grc{τῇ Βάαλ} (11:4), con artículo femenino: los lectores sustituían «Baal» por \heb{בֹּשֶׁת} (\emph{bóshet}), «vergüenza», en griego \grc{αἰσχύνη} (femenino). Pablo conocía la LXX y el hebreo (Stanley, 1992; Wagner, 2002).}\end{minipage}
\end{frame}

\begin{frame}{Las metáforas}
\begin{columns}[T,onlytextwidth]
\begin{column}{0.32\textwidth}\small\RaggedRight
\textbf{\textcolor{ladrillo}{Primicias y masa}} (11:16a)\par\smallskip
Nm 15:20: \heb{רֵאשִׁית עֲרִסֹתֵכֶם}; LXX \grc{ἀπαρχὴν φυράματος}. La porción ofrecida santifica la masa. Tratado \emph{Jalá} de la Mishná.
\end{column}
\begin{column}{0.33\textwidth}\small\RaggedRight
\textbf{\textcolor{ladrillo}{Raíz y ramas}} (11:16b)\par\smallskip
«Si la raíz es santa, también las ramas». La raíz: los padres (11:28), según la mayoría de los comentaristas (Moo, 1996).
\end{column}
\begin{column}{0.32\textwidth}\small\RaggedRight
\textbf{\textcolor{ladrillo}{El olivo}} (11:17--24)\par\smallskip
\heb{זַיִת רַעֲנָן}, «olivo verde», nombre de Judá (Jer 11:16); «su esplendor como el olivo» (Os 14:6 [14:7]); Sal 52:8 [52:10].
\end{column}
\end{columns}\vspace{8pt}
\begin{tcolorbox}[colback=piedra,colframe=cafe,boxrule=0.8pt,arc=3pt,fontupper=\small]
Un solo árbol, una sola raíz. Las ramas desgajadas siguen siendo «naturales» (\grc{κατὰ φύσιν}); los injertados siguen siendo «silvestres por naturaleza» (11:24). Ambos permanecen por la fe (11:20, 23).
\end{tcolorbox}
\end{frame}

\begin{frame}{La doxología en quiasmo}
\begin{columns}[c,onlytextwidth]
\begin{column}{0.50\textwidth}
\begin{cita}[fontupper=\scriptsize\itshape]¡Oh profundidad de las riquezas, y de la sabiduría y del conocimiento de Dios! [\ldots] ¿Quién conoció la mente del Señor? ¿O quién fue su consejero? ¿O quién le dio a él primero? [\ldots] Porque de él, por él y para él son todas las cosas.\fuente{(Ro 11:33--36)}\end{cita}
\scriptsize Is 40:13 (consuelo del exilio) y Job 41:11 [41:3]; \grc{ἀνεξιχνίαστοι}, cf. Job 5:9; 9:10 LXX.
\end{column}
\begin{column}{0.47\textwidth}
\centering
\begin{tikzpicture}[y=-0.62cm]
\foreach \i/\t/\x in {0/{a. riqueza}/0,1/{b. sabiduría}/0.5,2/{c. conocimiento}/1.0,3/{c'. ¿quién conoció la mente?}/1.0,4/{b'. ¿quién fue su consejero?}/0.5,5/{a'. ¿quién le dio primero?}/0}{
  \node[caja,anchor=west,inner sep=3pt] at (\x,\i) {\t};}
\node[font=\tiny,text=sepia,anchor=west] at (4.6,1) {v.\,33};
\node[font=\tiny,text=sepia,anchor=west] at (4.6,4) {vv.\,34--35};
\node[caja,fill=oro!25,anchor=west] at (0,6.4) {\grc{ἐξ αὐτοῦ καὶ δι᾿ αὐτοῦ καὶ εἰς αὐτόν} (v.\,36)};
\end{tikzpicture}
\end{column}
\end{columns}
\end{frame}

% ===================================================
\section{Filología}
\begin{frame}{Ro 11:25--27 en griego}
\begin{cita}[fontupper=\upshape\small]\grc{Οὐ γὰρ θέλω ὑμᾶς ἀγνοεῖν, ἀδελφοί, τὸ μυστήριον τοῦτο, ἵνα μὴ ἦτε [παρ᾿] ἑαυτοῖς φρόνιμοι, ὅτι πώρωσις ἀπὸ μέρους τῷ Ἰσραὴλ γέγονεν ἄχρι οὗ τὸ πλήρωμα τῶν ἐθνῶν εἰσέλθῃ καὶ οὕτως πᾶς Ἰσραὴλ σωθήσεται, καθὼς γέγραπται· ἥξει ἐκ Σιὼν ὁ ῥυόμενος, ἀποστρέψει ἀσεβείας ἀπὸ Ἰακώβ. καὶ αὕτη αὐτοῖς ἡ παρ᾿ ἐμοῦ διαθήκη, ὅταν ἀφέλωμαι τὰς ἁμαρτίας αὐτῶν.}\end{cita}
\begin{cita}[fontupper=\scriptsize\itshape]No quiero, hermanos, que ignoréis este misterio, para que no seáis sabios en vuestra propia opinión: que un endurecimiento en parte le ha sobrevenido a Israel hasta que haya entrado la plenitud de las naciones; y así todo Israel será salvo, como está escrito: «Vendrá de Sion el Libertador, apartará de Jacob las impiedades. Y este será mi pacto con ellos, cuando yo quite sus pecados».\fuente{(Ro 11:25--27, NA28; trad. propia)}\end{cita}
\precision{variantes menores: 11:6 (adición bizantina), 11:17 («de la raíz \emph{y} de la grosura»), 11:25 (\grc{παρ᾿}/\grc{ἐν}), 11:31 (segundo \grc{νῦν} entre corchetes; véase «Pasaje por pasaje»).}
\end{frame}

\begin{frame}{Elección y remanente}
\begin{columns}[T,onlytextwidth]
\begin{column}{0.46\textwidth}\small\RaggedRight
\begin{itemize}
\item \grc{ἀπώσατο} (11:1): eco de \grc{οὐκ ἀπώσεται κύριος τὸν λαὸν αὐτοῦ} (1~S 12:22; Sal 94:14 [LXX 93:14]).
\item \heb{כִּי לֹא־יִטֹּשׁ יְהוָה אֶת־עַמּוֹ}: «YHWH no desamparará a su pueblo» (1~S 12:22; Estudio 3).
\item \grc{προέγνω} (11:2): «conoció» = eligió (Am 3:2). La elección precede a la incredulidad.
\end{itemize}
\end{column}
\begin{column}{0.52\textwidth}
\centering
\begin{tikzpicture}[y=-1.15cm]
\node[caja,text width=4.6cm,align=center] (a) at (0,0) {\textbf{Gn 45:7}: \heb{שְׁאֵרִית} (José; Est. 1)};
\node[caja,text width=4.6cm,align=center] (b) at (0,1) {\textbf{1~R 19:18}: siete mil (reino del norte; Est. 5)};
\node[caja,text width=4.6cm,align=center] (c) at (0,2) {\textbf{Is 10:21}: \heb{שְׁאָר יָשׁוּב}, «un resto volverá»};
\node[caja,text width=4.6cm,align=center,fill=oro!25] (d) at (0,3) {\textbf{Ro 11:5}: \grc{λεῖμμα κατ᾿ ἐκλογὴν χάριτος}};
\foreach \u/\v in {a/b,b/c,c/d}{\draw[-{Latex[length=2mm]},oro!90!ladrillo,line width=1.2pt] (\u) -- (\v);}
\end{tikzpicture}\par\smallskip
{\scriptsize\itshape\color{sepia}\grc{γέγονεν}, perfecto: el remanente existe ahora y garantiza que la historia continúa.}
\end{column}
\end{columns}
\end{frame}

\begin{frame}{Endurecimiento, tropiezo y celos}
\begin{columns}[T,onlytextwidth]
\begin{column}{0.49\textwidth}\small\RaggedRight
\textbf{\textcolor{cafe}{\grc{πώρωσις}} (11:7, 25)}: de \grc{πῶρος}, callo óseo, piedra porosa: insensibilidad. Velo «hasta el día de hoy» (2~Co 3:14--16).\par\smallskip
\textbf{\textcolor{cafe}{\grc{πνεῦμα κατανύξεως}} (11:8)}: Is 29:10 \heb{רוּחַ תַּרְדֵּמָה}, el «sueño profundo» de Gn 2:21 y 15:12.\par\smallskip
\textbf{\textcolor{cafe}{\grc{ἔπταισαν ἵνα πέσωσιν;}} (11:11)}: tropezar (\grc{πταίω}) no es caer (\grc{πίπτω}). \grc{παράπτωμα}: «paso en falso».
\end{column}
\begin{column}{0.49\textwidth}
\centering{\heb{וַאֲנִי אַקְנִיאֵם בְּלֹא־עָם}}\par\smallskip
\begin{cita}[fontupper=\scriptsize\itshape]Yo también los provocaré a celos con lo que no es pueblo.\fuente{(Dt 32:21; LXX \grc{παραζηλώσω})}\end{cita}
\begin{cita}[fontupper=\scriptsize\itshape]Por si de algún modo provoco a celos a los de mi carne, y salvo a algunos de ellos.\fuente{(Ro 11:14)}\end{cita}
\scriptsize La misión a los gentiles tiene una finalidad que vuelve sobre Israel (Estudio 2).
\end{column}
\end{columns}
\end{frame}

\begin{frame}{Plenitud y aceptación}
\begin{columns}[T,onlytextwidth]
\begin{column}{0.49\textwidth}
\begin{cita}[fontupper=\scriptsize\itshape]Si su transgresión es la riqueza del mundo, y su disminución la riqueza de las naciones, ¡cuánto más su plenitud!\fuente{(Ro 11:12)}\end{cita}
\scriptsize
\begin{itemize}
\item \grc{ἥττημα}: «disminución», «derrota».
\item \grc{πλήρωμα αὐτῶν} (11:12) // \grc{πλήρωμα τῶν ἐθνῶν} (11:25): plenitud de Israel y de las naciones, en paralelo.
\end{itemize}
\precision{Gn 48:19 \heb{מְלֹא־הַגּוֹיִם} (Efraín) en la LXX es \grc{πλῆθος ἐθνῶν}, no \grc{πλήρωμα}; Pablo no lo cita.}
\end{column}
\begin{column}{0.49\textwidth}
\begin{cita}[fontupper=\scriptsize\itshape]Si su exclusión es la reconciliación del mundo, ¿qué será su aceptación sino vida de entre los muertos?\fuente{(Ro 11:15)}\end{cita}
\scriptsize
\begin{itemize}
\item \grc{ἀποβολή}: exclusión (temporal) por parte de Dios.
\item \grc{πρόσλημψις}: «acogida»; hápax del NT.
\item \grc{ζωὴ ἐκ νεκρῶν}: resurrección o bendición inmensa; trasfondo de Ez 37:11--14 (Estudio 5).
\item \grc{ἀπαρχή}/\grc{φύραμα} (11:16): \heb{רֵאשִׁית} y la masa.
\end{itemize}
\end{column}
\end{columns}
\end{frame}

\begin{frame}{El vocabulario del olivo}
\begin{columns}[T,onlytextwidth]
\begin{column}{0.49\textwidth}\small\RaggedRight
\begin{itemize}
\item \grc{ἀγριέλαιος}: olivo silvestre, acebuche.
\item \grc{καλλιέλαιος}: olivo cultivado, «hermoso».
\item \grc{ἐγκεντρίζω}: injertar por incisión; 6 veces en 11:17--24 (conteo propio).
\item \grc{ἐξεκλάσθησαν}: «fueron desgajadas».
\item \grc{συγκοινωνὸς τῆς ῥίζης τῆς πιότητος}: copartícipe de la savia de la raíz.
\item \grc{χρηστότης} y \grc{ἀποτομία} (11:22): bondad y «corte tajante».
\end{itemize}
\end{column}
\begin{column}{0.49\textwidth}
\begin{tcolorbox}[colback=piedra,colframe=cafe,boxrule=0.8pt,arc=3pt,fontupper=\scriptsize,title={\scriptsize ¿Sabía Pablo injertar?},colbacktitle=cafe,coltitle=piedra]
Lo común era injertar la púa cultivada en patrón silvestre. Pablo describe lo contrario y lo llama \grc{παρὰ φύσιν}, «contra la naturaleza» (11:24). Columela (\emph{De re rustica} 5.9.16) describe insertar una rama de acebuche en olivos viejos para revitalizarlos (Baxter y Ziesler, 1985). Debate abierto: práctica real conocida o inversión deliberada para subrayar la gracia.
\end{tcolorbox}
\end{column}
\end{columns}
\end{frame}

\begin{frame}{«Para que no seáis sabios en vuestra opinión»}
\begin{columns}[T,onlytextwidth]
\begin{column}{0.52\textwidth}\scriptsize\RaggedRight
\grc{ἵνα μὴ ἦτε [παρ᾿] ἑαυτοῖς φρόνιμοι} (Ro 11:25)\par\smallskip
\begin{itemize}
\item Eco de Pr 3:7 LXX: \grc{μὴ ἴσθι φρόνιμος παρὰ σεαυτῷ}; hebreo \heb{אַל־תְּהִי חָכָם בְּעֵינֶיךָ}, «no seas sabio a tus propios ojos».
\item Pablo lo repite en Ro 12:16.
\item RVR1960: «para que no seáis \textbf{arrogantes} en cuanto a vosotros mismos»; NVI: «presuntuosos».
\end{itemize}
\end{column}
\begin{column}{0.46\textwidth}\scriptsize\RaggedRight
\textbf{\textcolor{cafe}{\grc{μυστήριον}}}: 28 veces en el NT, 21 en Pablo (conteo propio).\par\smallskip
Agustín: \emph{Novum in Vetere latet, Vetus in Novo patet} (\emph{Quaest. in Hept.} 2.73).\par\medskip
\begin{tcolorbox}[colback=piedra,colframe=ladrillo,boxrule=0.8pt,arc=2pt,fontupper=\scriptsize]
Ignorar el plan de Dios con Israel no es un error inocente: \textbf{alimenta la soberbia} gentil (11:18, 20).
\end{tcolorbox}
\end{column}
\end{columns}
\end{frame}

\begin{frame}{Misterio, «en parte», «hasta que»}
\begin{columns}[T,onlytextwidth]
\begin{column}{0.33\textwidth}\small\RaggedRight
\textbf{\textcolor{ladrillo}{\grc{μυστήριον}}}\par\smallskip
Designio antes oculto, ahora revelado (Ro 16:25--26). Trasfondo: arameo \heb{רָז} (\emph{raz}) de Dn 2:18--30; Qumrán (1QpHab 7). Lo nuevo no es la restauración, sino \destacar{el orden}.
\end{column}
\begin{column}{0.31\textwidth}\small\RaggedRight
\textbf{\textcolor{ladrillo}{\grc{ἀπὸ μέρους}}}\par\smallskip
«En parte»: endurecimiento parcial o sobre una parte de Israel. Mismo resultado: el remanente no está endurecido (11:5, 7).
\end{column}
\begin{column}{0.33\textwidth}\small\RaggedRight
\textbf{\textcolor{ladrillo}{\grc{ἄχρι οὗ}}}\par\smallskip
«Hasta que» + subjuntivo: término de una situación.
\begin{itemize}
\item Lc 21:24: Jerusalén hollada \grc{ἄχρι οὗ} se cumplan los tiempos de los gentiles.
\item 1~Co 11:26; 15:25.
\end{itemize}
\grc{εἰσέλθῃ}: verbo de «entrar en el reino».
\end{column}
\end{columns}
\end{frame}

\begin{frame}{«Y así»: tres lecturas}
\centering\scriptsize\renewcommand{\arraystretch}{1.25}
\begin{tabular}{>{\RaggedRight\arraybackslash}p{2.3cm}>{\RaggedRight\arraybackslash}p{4.3cm}>{\RaggedRight\arraybackslash}p{4.4cm}}
\toprule
\textbf{\grc{καὶ οὕτως}} & \textbf{Sentido} & \textbf{Defensores}\\\midrule
Modal: «y de esta manera» & Mediante el proceso descrito: endurecimiento, entrada de las naciones, celos & Cranfield (1979); Moo (1996); Schreiner (1998); Fitzmyer (1993)\\
Temporal: «y entonces» & Después de la entrada de las naciones & Van der Horst (2000); RVR: «y luego»\\
Correlativa: «así, como está escrito» & \grc{οὕτως} anticipa \grc{καθὼς γέγραπται} & Matiz señalado por Dunn (1988)\\
\bottomrule
\end{tabular}\par\vspace{8pt}
\begin{tcolorbox}[colback=piedra,colframe=ladrillo,boxrule=1pt,arc=3pt,fontupper=\small,width=0.92\linewidth]
Las tres coinciden: la salvación de «todo Israel» \textbf{sigue} al endurecimiento parcial y a la entrada de la plenitud de las naciones. En 11:25, a un versículo, «Israel» se distingue de «las naciones».
\end{tcolorbox}
\end{frame}

\begin{frame}{\grc{πᾶς Ἰσραήλ}: «todo Israel»}
\begin{columns}[T,onlytextwidth]
\begin{column}{0.49\textwidth}\small
\textbf{\textcolor{cafe}{Uso colectivo en el AT}} (\heb{כָּל־יִשְׂרָאֵל})
\begin{itemize}
\item 1~S 25:1: todo Israel llora a Samuel.
\item 1~R 12:1: todo Israel en Siquem (Est. 4).
\item 2~Cr 12:1: «todo Israel» = el reino de Judá.
\item Dn 9:11: «todo Israel traspasó tu ley».
\item Neh 12:47: todo Israel sostiene a los cantores.
\end{itemize}
\scriptsize Nunca cada individuo; nunca no israelitas.
\end{column}
\begin{column}{0.49\textwidth}\small
\textbf{\textcolor{cafe}{Mishná, Sanedrín 10:1}}\par\smallskip
\centering{\heb{כל ישראל יש להם חלק לעולם הבא}}\par\smallskip
\raggedright «Todo Israel tiene parte en el mundo venidero» \ldots{} \textbf{y a continuación enumera quiénes no la tienen}.\par\medskip
\grc{Ἰσραήλ}: 11 veces en Ro 9--11, siempre el pueblo étnico (9:6b: un Israel \emph{dentro} de Israel).\par\smallskip
\grc{σωθήσεται}: futuro pasivo; obra de Dios; salvación por la fe (10:9--13).
\end{column}
\end{columns}
\end{frame}

\begin{frame}{Is 59:20: TM, LXX y Pablo}
\centering\scriptsize\renewcommand{\arraystretch}{1.35}
\begin{tabular}{>{\RaggedRight\arraybackslash}p{1.6cm}>{\RaggedRight\arraybackslash}p{4.3cm}>{\RaggedRight\arraybackslash}p{5.2cm}}
\toprule
\textbf{Versión} & \textbf{Texto} & \textbf{Traducción}\\\midrule
TM & \heb{וּבָא לְצִיּוֹן גּוֹאֵל וּלְשָׁבֵי פֶשַׁע בְּיַעֲקֹב} & «Vendrá \textbf{a} Sion un Redentor, y a los que se vuelven de la rebelión en Jacob»\\
LXX & \grc{ἥξει ἕνεκεν Σιων ὁ ῥυόμενος καὶ ἀποστρέψει ἀσεβείας ἀπὸ Ιακωβ} & «Vendrá \textbf{por causa de} Sion el que libra, y apartará las impiedades de Jacob»\\
Pablo & \grc{ἥξει ἐκ Σιὼν ὁ ῥυόμενος, ἀποστρέψει ἀσεβείας ἀπὸ Ἰακώβ} & «Vendrá \textbf{de} Sion el Libertador, apartará de Jacob las impiedades»\\
\bottomrule
\end{tabular}\par\vspace{6pt}
\begin{minipage}{0.95\linewidth}\scriptsize
\textbf{\grc{ἐκ Σιών}}: ecos de Sal 14:7 e Is 2:3; la Sion celestial y la segunda venida (Cranfield, 1979; Moo, 1996; cf. 1~Ts 1:10, \grc{τὸν ῥυόμενον}), o la misión desde Jerusalén (Wright, 2013).\par\smallskip
\textbf{Pacto}: Is 59:21 + 27:9 LXX («cuando yo quite sus pecados») evocan Jer 31:33--34 [31:32--33]: el nuevo pacto con Israel y Judá (Estudio 5).
\end{minipage}
\end{frame}

\begin{frame}{El Redentor de Sion}
\fotolado{panel_is59.jpg}{Elaboración propia: Is 59:20--21a en texto hebreo consonántico (BHS).}
\begin{columns}[c,onlytextwidth]
\begin{column}{0.57\textwidth}
\centering{\large\heb{וּבָא לְצִיּוֹן גּוֹאֵל}}\par\smallskip
\begin{cita}Y vendrá a Sion el Redentor, y a los que se vuelven de la rebelión en Jacob, dice YHWH. Y en cuanto a mí, este es mi pacto con ellos.\fuente{(Is 59:20--21)}\end{cita}
\scriptsize
\begin{itemize}
\item \heb{גּוֹאֵל} (\emph{go'el}): el pariente que rescata (Rut; Is 41:14).
\item El TM pone el arrepentimiento antes; la LXX y Pablo, la acción del Libertador.
\item Estos versículos abren la oración diaria \emph{U-va le-Tsiyyón} de la liturgia judía.
\end{itemize}
\end{column}
\begin{column}{0.41\textwidth}\end{column}
\end{columns}
\end{frame}

\begin{frame}{«Y así todo Israel»}
\fotolado{panel_ro26.jpg}{Elaboración propia: Ro 11:25b--26 en escritura uncial (texto NA28).}
\begin{columns}[c,onlytextwidth]
\begin{column}{0.57\textwidth}
\begin{cita}Un endurecimiento en parte le ha sobrevenido a Israel hasta que haya entrado la plenitud de las naciones; y así todo Israel será salvo.\fuente{(Ro 11:25b--26a)}\end{cita}
\scriptsize
\begin{itemize}
\item \grc{ἀμεταμέλητα} (11:29): «sin arrepentimiento»; Nm 23:19.
\item \grc{ἐχθροί … ἀγαπητοί} (11:28): enemigos y amados \emph{a la vez}.
\item \grc{συνέκλεισεν} (11:32): «encerró juntos», como la red (Lc 5:6; Gá 3:22).
\item \grc{βάθος} (11:33): la última palabra es «profundidad».
\end{itemize}
\end{column}
\begin{column}{0.41\textwidth}\end{column}
\end{columns}
\end{frame}

% ===================================================
\section{Teología}
\begin{frame}{La fidelidad de Dios}
\begin{columns}[T,onlytextwidth]
\begin{column}{0.49\textwidth}\small\RaggedRight
Si Dios abandonó a Israel, ¿qué valor tiene Ro 8:39, «nada nos podrá separar»?\par\medskip
\destacar{Ro 9:6}: «no es que la palabra de Dios haya fallado».\par\medskip
Ro 11 completa la prueba: la elección precede a la historia (\grc{προέγνω}) y es irrevocable (\grc{ἀμεταμέλητα}).
\end{column}
\begin{column}{0.49\textwidth}\small\RaggedRight
\textbf{\textcolor{cafe}{El endurecimiento tiene tres límites}}
\begin{enumerate}
\item \textbf{Parcial} (\grc{ἀπὸ μέρους}): existe el remanente.
\item \textbf{Temporal} (\grc{ἄχρι οὗ}): tiene término.
\item \textbf{Con propósito}: salvación de las naciones y celos de Israel (11:11).
\end{enumerate}
\scriptsize Razones según la Clase 7: soberanía (9:18), ignorancia (10:4), justicia propia (10:3), temporalidad (11:26).
\end{column}
\end{columns}
\end{frame}

\begin{frame}{El olivo}
\centering
\begin{tikzpicture}[x=0.95cm,y=0.68cm]
\fill[cafe!75] (-0.35,0) -- (-0.28,3.1) -- (0.28,3.1) -- (0.35,0) -- cycle;
\foreach \a in {-55,-28,0,28,55}{\draw[cafe,line width=1.4pt] (0,0) -- ++(-90+\a:1.0);}
\node[font=\scriptsize\bfseries,text=ladrillo] at (0,-1.55) {Raíz santa: los padres (11:16, 28)};
\foreach \x/\y in {-1.7/4.5, 0/5.0, 1.7/4.5}{\draw[oro!75!cafe,line width=2.4pt] (0,2.9) -- (\x,\y); \fill[oro!85!cafe] (\x,\y) circle (5pt);}
\node[font=\tiny,text=cafeosc,text width=3.2cm,align=center] at (0,5.7) {Ramas naturales que permanecen: el remanente (11:5)};
\draw[ladrillo,line width=2.4pt] (0,2.2) -- (-2.4,3.0); \fill[ladrillo] (-2.4,3.0) circle (4.5pt);
\draw[ladrillo,line width=2.4pt] (0,2.0) -- (2.4,2.9); \fill[ladrillo] (2.4,2.9) circle (4.5pt);
\node[font=\tiny,text=ladrillo,text width=2.8cm,align=center] at (-4.7,3.3) {\textbf{Ramas silvestres injertadas} por la fe (11:17, 20)};
\draw[oro!75!cafe,line width=2.4pt,dashed] (-2.6,1.3) -- (-3.5,0.7); \fill[oro!45!piedra,draw=oro!75!cafe] (-3.5,0.7) circle (4.5pt);
\node[font=\tiny,text=sepia,text width=2.8cm,align=center] at (-4.7,1.4) {Ramas naturales desgajadas por incredulidad (11:17, 20)};
\draw[oro!75!cafe,line width=2.4pt,dashed] (2.9,0.9) -- (3.8,0.4); \fill[oro!45!piedra,draw=oro!75!cafe] (3.8,0.4) circle (4.5pt);
\draw[oro!75!cafe,line width=2.4pt] (0,1.5) -- (2.2,1.8); \fill[oro!85!cafe] (2.2,1.8) circle (4.5pt);
\draw[-{Latex[length=2mm]},cafeosc,line width=0.9pt,dashed] (3.7,0.62) to[bend right=25] (2.45,1.75);
\node[font=\tiny,text=cafeosc,text width=3.0cm,align=center] at (5.0,2.6) {\textbf{Reinjertadas} «si no permanecen en incredulidad» (11:23--24)};
\end{tikzpicture}\par\vspace{2pt}
{\scriptsize Injertados «entre» las ramas (\grc{ἐν αὐτοῖς}): \destacar{inclusión sin reemplazo}. «No sustentas tú a la raíz, sino la raíz a ti» (11:18).}\par
{\fontsize{5}{6}\selectfont\color{sepia}\raggedleft\textit{Nota.} Elaboración propia: esquema del olivo de Ro 11:16--24.\par}
\end{frame}

\begin{frame}{El misterio: la secuencia}
\centering
\begin{tikzpicture}[node distance=0.5cm]
\node[caja,text width=3.0cm,align=center,minimum height=2.0cm] (a) {\small\textbf{\textcolor{ladrillo}{Endurecimiento en parte}}\\[3pt]\grc{πώρωσις}\\\grc{ἀπὸ μέρους}\\[2pt]{\tiny remanente creyente + mayoría endurecida}};
\node[caja,text width=3.0cm,align=center,minimum height=2.0cm,right=of a] (b) {\small\textbf{\textcolor{ladrillo}{Entra la plenitud de las naciones}}\\[3pt]\grc{ἄχρι οὗ τὸ πλήρωμα}\\\grc{τῶν ἐθνῶν εἰσέλθῃ}};
\node[caja,text width=3.0cm,align=center,minimum height=2.0cm,right=of b,fill=oro!25] (c) {\small\textbf{\textcolor{ladrillo}{Y así todo Israel será salvo}}\\[3pt]\grc{καὶ οὕτως πᾶς}\\\grc{Ἰσραὴλ σωθήσεται}};
\draw[-{Latex[length=2.5mm]},oro!90!ladrillo,line width=1.4pt] (a) -- (b);
\draw[-{Latex[length=2.5mm]},oro!90!ladrillo,line width=1.4pt] (b) -- (c);
\end{tikzpicture}\par\vspace{8pt}
\begin{cita}[fontupper=\scriptsize\itshape]Vendrá de Sion el Libertador, apartará de Jacob las impiedades. Y este será mi pacto con ellos, cuando yo quite sus pecados.\fuente{(Ro 11:26b--27; Is 59:20--21; 27:9; Jer 31:33--34)}\end{cita}
{\scriptsize Los profetas: Israel restaurado y las naciones acuden (Is 2:2--4; Zac 8:23). \destacar{El misterio invierte el orden}: primero entran las naciones.}\par
{\fontsize{5}{6}\selectfont\color{sepia}\raggedleft\textit{Nota.} Elaboración propia: la secuencia de Ro 11:25--27.\par}
\end{frame}

\begin{frame}{¿Por qué tropezó Israel?}
\begin{columns}[T,onlytextwidth]
\begin{column}{0.49\textwidth}\scriptsize\RaggedRight
\textbf{\textcolor{cafe}{La piedra de tropiezo}}
\begin{itemize}
\item «Maldito por Dios es el colgado» (Dt 21:23; Gá 3:13).
\item Cristo crucificado: \grc{σκάνδαλον} para los judíos (1~Co 1:23).
\item La «piedra de tropiezo» de Is 8:14; 28:16 es el Mesías (Ro 9:32--33).
\item Esperaban un rey libertador de Roma; recibieron un siervo sufriente (Is 53).
\end{itemize}
\end{column}
\begin{column}{0.49\textwidth}\scriptsize\RaggedRight
\textbf{\textcolor{cafe}{Los gentiles también estaban endurecidos}}
\begin{itemize}
\item «Dios los \textbf{entregó}» (\grc{παρέδωκεν}): Ro 1:24, 26, 28.
\item «Sin Cristo, ajenos a los pactos, sin esperanza» (Ef 2:12).
\item Ef 4:18: «por la \textbf{dureza} de su corazón», \grc{διὰ τὴν πώρωσιν}: \textbf{la misma palabra} de Ro 11:25.
\end{itemize}
\begin{tcolorbox}[colback=piedra,colframe=ladrillo,boxrule=0.8pt,arc=2pt,fontupper=\scriptsize]
El \emph{pṓrōsis} de Israel no es peor que el de las naciones: «Dios encerró a todos» (11:32).
\end{tcolorbox}
\end{column}
\end{columns}
\end{frame}

\begin{frame}{Amados por los padres}
\fotolado{panel_jer31.jpg}{Elaboración propia: Jer 31:35a, 36 en texto hebreo consonántico (BHS).}
\begin{columns}[c,onlytextwidth]
\begin{column}{0.57\textwidth}
\begin{cita}[fontupper=\scriptsize\itshape]En cuanto al evangelio, enemigos por causa de vosotros; en cuanto a la elección, amados por causa de los padres.\fuente{(Ro 11:28)}\end{cita}
\centering{\heb{גַּם זֶרַע יִשְׂרָאֵל יִשְׁבְּתוּ מִהְיוֹת גּוֹי}}\par\smallskip
\begin{cita}[fontupper=\scriptsize\itshape]Si faltaren estas leyes delante de mí, dice YHWH, también la descendencia de Israel dejará de ser nación delante de mí todos los días.\fuente{(Jer 31:36 [31:35])}\end{cita}
\scriptsize El pacto que perdona (Jer 31:31--34) es el pacto que conserva al pueblo (31:35--37).
\end{column}
\begin{column}{0.41\textwidth}\end{column}
\end{columns}
\end{frame}

\begin{frame}{«Por amor de Jonatán tu padre»}
\begin{columns}[T,onlytextwidth]
\begin{column}{0.52\textwidth}\scriptsize\RaggedRight
\textbf{\textcolor{cafe}{La base: Dt 7:7--8}}\par
«No por ser vosotros más que todos los pueblos [\ldots] sino por cuanto YHWH os amó, y quiso guardar el juramento que juró a \textbf{vuestros padres}».\par\medskip
\textbf{\textcolor{cafe}{La imagen: Mefiboset}} (Estudio 4)
\begin{itemize}
\item Pacto de David con Jonatán (1~S 20:14--17, 42).
\item «Haré contigo misericordia (\heb{חֶסֶד}) \textbf{por amor de Jonatán tu padre}» (2~S 9:7).
\item Lisiado, sin méritos, come siempre a la mesa del rey.
\end{itemize}
\end{column}
\begin{column}{0.46\textwidth}\scriptsize\RaggedRight
\textbf{\textcolor{cafe}{Enemigos y amados a la vez}}\par
El hijo pródigo (Lc 15:11--32): actúa como enemigo de la casa mientras el padre lo ama y lo espera.\par\smallskip
\precision{la parábola habla de un individuo; Ro 11:28, de un pueblo y una elección colectiva.}\par\medskip
\begin{tcolorbox}[colback=piedra,colframe=ladrillo,boxrule=0.8pt,arc=2pt,fontupper=\scriptsize]
Israel es amado \textbf{por causa de los padres}: por el pacto, no por la conducta de cada generación (Ro 11:28--29; Nm 23:19).
\end{tcolorbox}
\end{column}
\end{columns}
\end{frame}

\begin{frame}{Misericordia para todos}
\begin{columns}[c,onlytextwidth]
\begin{column}{0.52\textwidth}
\begin{cita}Porque Dios encerró a todos en desobediencia, para tener misericordia de todos.\fuente{(Ro 11:32)}\end{cita}
\small «Todos» (\grc{τοὺς πάντας}): los dos grupos del contexto, gentiles e Israel.\par\medskip
Ningún grupo reclama mérito: ambos se salvan por la misma misericordia.
\end{column}
\begin{column}{0.45\textwidth}
\centering
\begin{tikzpicture}[node distance=0.6cm]
\node[caja,text width=3.6cm,align=center] (g) {\textbf{Gentiles}: desobedientes $\to$ misericordia por la desobediencia de Israel (11:30)};
\node[caja,text width=3.6cm,align=center,below=of g] (i) {\textbf{Israel}: desobediente $\to$ misericordia por la misericordia a los gentiles (11:31)};
\draw[-{Latex[length=2mm]},oro!90!ladrillo,line width=1.2pt] (g.east) to[bend left=50] (i.east);
\draw[-{Latex[length=2mm]},oro!90!ladrillo,line width=1.2pt] (i.west) to[bend left=50] (g.west);
\end{tikzpicture}
\end{column}
\end{columns}
\end{frame}

\begin{frame}{Nación santa en Romanos 11}
\centering
\begin{tikzpicture}
 \foreach \x/\n/\t/\d/\r in {0/1/{Santidad}/{«Si las primicias son santas, también la masa; si la raíz es santa, también las ramas».}/{Ro 11:16; Éx 19:6},
   4.2/2/{Mediación}/{Por el tropiezo de Israel «vino la salvación a las naciones»; el remanente apostólico llevó el evangelio. Pablo, «sacerdote del evangelio».}/{Ro 11:11; 15:16},
   8.4/3/{Plenitud}/{Si el tropiezo fue riqueza, la plenitud será «vida de entre los muertos». «Seréis llamados sacerdotes de YHWH».}/{Ro 11:12, 15; Is 61:6}}{
  \node[caja,anchor=north] (b\n) at (\x,0) {\parbox[t][3.0cm][t]{3.55cm}{\RaggedRight
    {\small\textbf{\textcolor{ladrillo}{\t}}}\\[5pt]\scriptsize \d\par\vfill{\tiny\color{sepia}(\r)}}};
  \node[circle,fill=cafe,text=piedra,font=\scriptsize\bfseries,inner sep=1.6pt] at (b\n.north west) {\n};}
\end{tikzpicture}\par\vspace{6pt}
{\small\itshape\color{sepia}La vocación sacerdotal de Israel no está cancelada: espera su forma completa.}
\end{frame}

% ===================================================
\section{Pasaje por pasaje}
\newcommand{\tabpost}[6]{%
\centering\scriptsize\renewcommand{\arraystretch}{1.18}
\begin{tabular}{>{\RaggedRight\arraybackslash}p{2.3cm}>{\RaggedRight\arraybackslash}p{8.9cm}}
\toprule
\textbf{Postura} & \textbf{Cómo lee el pasaje}\\\midrule
1. Iglesia & #1\\
2. Remanente & #2\\
3. Todos & #3\\
4. Dos caminos & #4\\
5. Dos Casas & #5\\
6. Israel futuro & #6\\
\bottomrule
\end{tabular}}

\begin{frame}{Seis posturas, ocho pasajes}
\begin{columns}[T,onlytextwidth]
\begin{column}{0.52\textwidth}\scriptsize
\textbf{\textcolor{cafe}{Clave de las posturas}}
\begin{itemize}
\item \textbf{1.} La iglesia (Calvino; Wright).
\item \textbf{2.} Remanente de todas las épocas (Hendriksen; Robertson; \textbf{Sugel Michelén}).
\item \textbf{3.} Todos los judíos.
\item \textbf{4.} Dos caminos (Stendahl; Gaston; Gager).
\item \textbf{5.} Dos Casas (Koniuchowsky).
\item \textbf{6.} Israel étnico futuro que reconoce al Mesías (Murray; Cranfield; Moo; Schreiner).
\end{itemize}
\end{column}
\begin{column}{0.46\textwidth}\scriptsize
\textbf{\textcolor{cafe}{Michelén, según el resumen compartido}}
\begin{itemize}
\item «Todo Israel» = suma de los elegidos en la historia.
\item Apoyo clave: el «\textbf{ahora}» de 11:30--31.
\item La misericordia funciona como un \textbf{vaivén} entre judíos y gentiles.
\item Un solo olivo: una sola familia de Dios.
\end{itemize}
\precision{postura 2 con matiz de la 1. Se examinan sus apoyos: «ahora» y «así».}
\end{column}
\end{columns}
\end{frame}

\begin{frame}{Pasaje 1 (11:1--6): remanente presente}
\begin{cita}[fontupper=\scriptsize\itshape]\grc{οὕτως οὖν καὶ ἐν τῷ νῦν καιρῷ λεῖμμα κατ᾿ ἐκλογὴν χάριτος γέγονεν}: «así también en el tiempo presente ha llegado a existir un remanente según la elección de la gracia».\fuente{(Ro 11:5)}\end{cita}
\tabpost{El remanente, con los gentiles, es el verdadero Israel.}%
{El remanente de hoy, sumado en la historia, es «todo Israel».}%
{«No desechó» = ningún judío fuera; pero solo «la elección» alcanzó (11:7).}%
{El remanente que Pablo muestra cree en Jesús: él mismo.}%
{Elías y los siete mil eran del norte; Pablo los usa como modelo, no como Efraín oculto.}%
{El remanente es prenda y comienzo; se opone a la «plenitud» (11:12).}
\par\vspace{2pt}\precision{el «ahora» de Pablo (\grc{νῦν καιρός}) es el tiempo \emph{del remanente}. Variante 11:1: P46 y G leen «su heredad».}
\end{frame}

\begin{frame}{Pasaje 2 (11:7--10): elegidos y endurecidos}
\begin{cita}[fontupper=\scriptsize\itshape]«La elección lo alcanzó, y los demás fueron endurecidos» \ldots{} «encorva su espalda \grc{διὰ παντός}» (\heb{תָּמִיד}, «continuamente»).\fuente{(Ro 11:7, 10; Sal 69:23)}\end{cita}
\tabpost{«Los demás» quedan fuera; lo de Israel que permanece entra en la iglesia.}%
{Cada generación se divide en elegidos y endurecidos.}%
{No explica que «los demás» no alcanzaran lo que buscaban.}%
{El endurecimiento abriría paso a los gentiles sin que Israel deba creer; Pablo lo ve como velo que se quita (2~Co 3:14--16).}%
{Sin lectura propia: habla del Israel judío de Pablo.}%
{Endurecimiento real, parcial y temporal: «continuamente» mientras dure; 11:25 le pone límite.}
\end{frame}

\begin{frame}{Pasaje 3 (11:11--12): el vaivén aparece}
\begin{cita}[fontupper=\scriptsize\itshape]Por su transgresión vino la salvación a las naciones, \textbf{para provocarlos a celos}. Y si su transgresión es riqueza para el mundo, ¡\textbf{cuánto más} su plenitud!\fuente{(Ro 11:11--12)}\end{cita}
\tabpost{La «plenitud» sería su entrada en la iglesia; pero no sería un bien \emph{mayor}.}%
{El vaivén ya opera; la plenitud es la suma de los atraídos.}%
{Plenitud = cada judío; el «cuánto más» no lo exige.}%
{Los celos buscan que Israel reciba la misma salvación (11:14).}%
{Plenitud de las tribus; pero es la plenitud de «ellos», los que tropezaron.}%
{Vaivén presente que apunta a un resultado mayor: plenitud frente a remanente.}
\par\vspace{2pt}\precision{\emph{qal va-jómer}: de lo menor (\grc{ἥττημα}) a lo mayor (\grc{πλήρωμα}).}
\end{frame}

\begin{frame}{Pasaje 4 (11:13--16): «algunos» y «vida de entre los muertos»}
\begin{cita}[fontupper=\scriptsize\itshape]Por si de algún modo provoco a celos a los de mi carne y salvo a \textbf{algunos} de ellos \ldots{} ¿qué será su aceptación sino vida de entre los muertos?\fuente{(Ro 11:14--15)}\end{cita}
\centering\scriptsize
\begin{tikzpicture}[x=2.85cm]
\foreach \i/\g/\v/\t in {0/λεῖμμα/11:5/remanente,1/τινάς/11:14/algunos,2/πλήρωμα/11:12/plenitud,3/πᾶς Ἰσραήλ/11:26/todo Israel}
  \node[caja,minimum width=2.25cm,text width=2.05cm,align=center] (n\i) at (\i,0) {\grc{\g}\\«\t» (\v)};
\draw[-{Stealth},oro,line width=1.2pt] (n0)--(n1); \draw[-{Stealth},oro,line width=1.2pt] (n1)--(n2); \draw[-{Stealth},oro,line width=1.2pt] (n2)--(n3);
\draw[decorate,decoration={brace,mirror,amplitude=4pt},cafe] ([yshift=-3pt]n0.south west)--([yshift=-3pt]n1.south east) node[midway,below=5pt]{«ya»: el presente de Pablo};
\draw[decorate,decoration={brace,mirror,amplitude=4pt},ladrillo] ([yshift=-3pt]n2.south west)--([yshift=-3pt]n3.south east) node[midway,below=5pt]{«todavía no»};
\end{tikzpicture}\par\vspace{4pt}
\raggedright\scriptsize
\textbf{2 (Michelén)}: el ministerio de Pablo prueba que el proceso es presente. \textbf{6}: es presente, pero salva a «algunos»; la «aceptación» será de otro orden (Ez 37:11--14). \textbf{3}: «algunos» la contradice.
\end{frame}

\begin{frame}{Pasaje 5 (11:17--24): un solo olivo}
\begin{cita}[fontupper=\scriptsize\itshape]¡Cuánto más estos, las ramas \textbf{naturales} (\grc{κατὰ φύσιν}), \textbf{serán injertados} en \textbf{su propio} olivo (\grc{τῇ ἰδίᾳ ἐλαίᾳ})!\fuente{(Ro 11:24)}\end{cita}
\tabpost{Un solo pueblo; las distinciones de origen se superan.}%
{Un solo olivo, una sola familia; reinjerto cada vez que un judío cree.}%
{Choca con «si no permanecen en incredulidad» (11:23).}%
{Choca con la misma condición de fe.}%
{Choca con «olivo silvestre por naturaleza»: Efraín serían ramas naturales.}%
{Un solo olivo, pero «su propio olivo»; futuro pasivo: reinjerto por venir.}
\par\vspace{2pt}\precision{Michelén acierta: no hay dos árboles. Pero el texto mantiene ramas naturales y silvestres.}
\end{frame}

\begin{frame}{«Y así» (\grc{καὶ οὕτως}): cierre de una secuencia}
\begin{columns}[T,onlytextwidth]
\begin{column}{0.53\textwidth}\scriptsize
\textbf{\textcolor{cafe}{\grc{καὶ οὕτως} al final de una serie de pasos}}
\begin{itemize}
\item Hch 27:44: \grc{καὶ οὕτως ἐγένετο πάντας διασωθῆναι}, «y así \textbf{todos} se \textbf{salvaron}».
\item 1~Ts 4:16--17: primero \ldots{} luego \ldots{} \grc{καὶ οὕτως} «estaremos siempre con el Señor».
\item Ro 5:12: «y así la muerte pasó a \textbf{todos}».
\item Hch 7:8; 28:14: «y así engendró», «y así llegamos».
\end{itemize}
Modo \emph{y} culminación de una secuencia. Van der Horst (2000): «solo entonces»; RVR: «y luego».
\end{column}
\begin{column}{0.45\textwidth}\scriptsize
\textbf{\textcolor{cafe}{Hacia atrás}}: resume el vaivén de 11:11--25 \ldots{} \textbf{incluido el «hasta que»}.\par\smallskip
\textbf{\textcolor{cafe}{Hacia adelante}}: \grc{οὕτως … καθὼς γέγραπται}, «así \ldots{} como está escrito» (cf. Lc 24:24): el modo lo define la cita: el Libertador quita la impiedad de \textbf{Jacob}.\par\smallskip
\begin{tcolorbox}[colback=piedra,colframe=ladrillo,boxrule=0.8pt,arc=2pt,fontupper=\scriptsize]
Mismo sujeto: «endurecimiento a \textbf{Israel} \ldots{} y así todo \textbf{Israel}». El «así» no puede cambiar quién es Israel.
\end{tcolorbox}
\end{column}
\end{columns}
\end{frame}

\begin{frame}{«Y así» según cada postura}
\centering\scriptsize\renewcommand{\arraystretch}{1.2}
\begin{tabular}{>{\RaggedRight\arraybackslash}p{2.2cm}>{\RaggedRight\arraybackslash}p{4.3cm}>{\RaggedRight\arraybackslash}p{4.6cm}}
\toprule
\textbf{Postura} & \textbf{«Y así» significa\ldots} & \textbf{Observación}\\\midrule
1. Iglesia & Por la entrada de gentiles y judíos creyentes & Cambia el sentido de «Israel» entre 11:25 y 26\\
2. Remanente & Por el vaivén continuo, generación tras generación & Buen antecedente; omite el término («hasta que»)\\
3. Todos & «Y así, sin más» & Ignora endurecimiento y fe\\
4. Dos caminos & Por el pacto, sin el evangelio & El vaivén es el del evangelio (11:11, 14)\\
5. Dos Casas & Al entrar los gentiles-Efraín & 11:24 lo excluye\\
6. Israel futuro & Por el vaivén que culmina al entrar la plenitud, como dice la Escritura & Mismo «Israel», «hasta que» y cita\\
\bottomrule
\end{tabular}
\end{frame}

\begin{frame}{¿«A Sion», «por Sion» o «de Sion»?}
\centering\scriptsize\renewcommand{\arraystretch}{1.2}
\begin{tabular}{>{\RaggedRight\arraybackslash}p{2.4cm}>{\RaggedRight\arraybackslash}p{2.4cm}>{\RaggedRight\arraybackslash}p{6.2cm}}
\toprule
\textbf{Testigo} & \textbf{Texto} & \textbf{Sentido}\\\midrule
Texto Masorético & \heb{וּבָא לְצִיּוֹן גּוֹאֵל} & «Vendrá \textbf{a} Sion un Redentor»: \heb{לְ} = «a, hacia, para»\\
1QIsa\textsuperscript{a} (Qumrán) & \heb{אל ציון} & «Hacia Sion»\\
Targum & \heb{לציון} & «A Sion» + volver a la Torá\\
Septuaginta & \grc{ἕνεκεν Σιων} & «\textbf{Por causa de} Sion»\\
Áquila, Símaco & \grc{τῇ Σιών} & «Para Sion»\\
Pablo (Ro 11:26) & \grc{ἐκ Σιών} & «\textbf{Desde} Sion»\\
\bottomrule
\end{tabular}\par\vspace{5pt}
\raggedright\scriptsize
\textbf{En español}: RVR traduce Is 59:20 «vendrá el Redentor \textbf{a} Sion» y Ro 11:26 «vendrá \textbf{de} Sion el Libertador». Cada versículo traduce fielmente su texto.\par\smallskip
\heb{גּוֹאֵל} (\emph{go'el}): el pariente que rescata (Lv 25:25; Rut 4) $\to$ \grc{ὁ ῥυόμενος}: Jesús, «que nos libra» y viene de los cielos (1~Ts 1:10).
\end{frame}

\begin{frame}{¿Por qué Pablo escribe «de Sion»?}
\begin{columns}[T,onlytextwidth]
\begin{column}{0.5\textwidth}\scriptsize
\begin{enumerate}
\item \textbf{Forma previa}: colección de textos de sinagoga o judeocreyente (Stanley, 1992).
\item \textbf{Combinación}: Is 2:3 / Mi 4:2, «de Sion saldrá la ley» (Bruno, 2008); Sal 14:7 [LXX 13:7], «¡quién dará desde Sion la salvación de Israel!» (Schreiner, 1998).
\item \textbf{Sion celestial}: segunda venida (Cranfield; Moo; Schreiner; cf. Gá 4:26; Heb 12:22).
\item \textbf{Sion terrenal}: primera venida y misión (Wright, 2013; Bruno, 2008).
\end{enumerate}
\end{column}
\begin{column}{0.48\textwidth}\scriptsize
\begin{tcolorbox}[colback=piedra,colframe=ladrillo,boxrule=0.8pt,arc=2pt,fontupper=\scriptsize,title={Evaluación},coltitle=piedra,colbacktitle=ladrillo]
Viene \textbf{a} Sion (TM), \textbf{por causa de} Sion (LXX) y actúa \textbf{desde} Sion (Pablo): un mismo acontecimiento.\par\smallskip
A favor del futuro: \grc{ἥξει} en futuro, como \grc{σωθήσεται}; el objeto es \textbf{Jacob}; Is 27:9 + Jer 31:34: perdón colectivo.
\end{tcolorbox}
\precision{el TM pone el arrepentimiento en Jacob; LXX y Pablo, la acción en el Libertador. Zac 12:10 une ambos.}
\end{column}
\end{columns}
\end{frame}

\begin{frame}{Pasaje 7 (11:30--31): la estructura del vaivén}
\centering\scriptsize\renewcommand{\arraystretch}{1.25}
\begin{tabular}{>{\RaggedRight\arraybackslash}p{1.8cm}>{\RaggedRight\arraybackslash}p{4.5cm}>{\RaggedRight\arraybackslash}p{4.7cm}}
\toprule
 & \textbf{Gentiles (11:30)} & \textbf{Israel (11:31)}\\\midrule
Desobediencia & \grc{ὑμεῖς} \textcolor{ladrillo}{\grc{ποτε}} \grc{ἠπειθήσατε}: «\emph{en otro tiempo} desobedecisteis» & \grc{οὗτοι} \textcolor{ladrillo}{\grc{νῦν}} \grc{ἠπείθησαν}: «\emph{ahora} han desobedecido»\\
Misericordia & \textcolor{ladrillo}{\grc{νῦν}} \grc{δὲ ἠλεήθητε}: «\emph{ahora} recibisteis misericordia» & \grc{ἵνα καὶ αὐτοὶ} \textcolor{ladrillo}{\grc{[νῦν]}} \grc{ἐλεηθῶσιν}: «para que [\emph{ahora}] la reciban»\\
Ocasión & «por la desobediencia de ellos» & «por la misericordia mostrada a vosotros»\\
\bottomrule
\end{tabular}\par\vspace{6pt}
\begin{tcolorbox}[colback=piedra,colframe=ladrillo,boxrule=1pt,arc=3pt,fontupper=\small,width=0.95\linewidth]
\textbf{El vaivén está establecido} sin depender del tercer «ahora»: 11:11, 12, 14, 15, 28, 30--32. La misericordia a las naciones es el medio de la misericordia a Israel. \textbf{En esto Michelén tiene razón.}
\end{tcolorbox}
\end{frame}

\begin{frame}{Los tres «ahora»: dos seguros, uno incierto}
\centering\scriptsize\renewcommand{\arraystretch}{1.2}
\begin{tabular}{>{\RaggedRight\arraybackslash}p{2.0cm}>{\RaggedRight\arraybackslash}p{6.3cm}>{\RaggedRight\arraybackslash}p{2.8cm}}
\toprule
\textbf{11:31b} & \textbf{Testigos} & \textbf{Familia}\\\midrule
\grc{νῦν}, «ahora» & Sinaítico (\heb{א}, s.~\textsc{iv}); Vaticano (B, s.~\textsc{iv}); Claromontano original (D*, s.~\textsc{vi}); 1506; copto bohaírico & Alejandrina y occidental\\
\grc{ὕστερον}, «después» & Minúsculos 33, 365; copto sahídico & Corrección: se leía como futuro\\
Sin adverbio & \textbf{P46 (c.~200)}; Alejandrino (A, s.~\textsc{v}); D\textsuperscript{2}; F, G (s.~\textsc{ix}); \grc{Ψ}; 1739; 1881; texto bizantino; latín & El más antiguo y la mayoría\\
\bottomrule
\end{tabular}\par\vspace{4pt}
\raggedright\scriptsize
\textbf{Ediciones}: NA28/UBS5 \grc{[νῦν]} entre corchetes; Westcott-Hort y Tischendorf lo incluyen; texto bizantino y TR de Scrivener lo omiten.\par
\textbf{Español}: Reina-Valera (1569, 1909, 1960): «ellos también alcancen misericordia», \textbf{sin «ahora»}; LBLA: «también a ellos \textbf{ahora}».\par\smallskip
{\tiny\color{sepia}Testigos según la nota textual de la NET Bible a Ro 11:31. En Pablo, D es el Códice Claromontano.}
\end{frame}

\begin{frame}{¿Cómo se decide el tercer «ahora»?}
\begin{columns}[T,onlytextwidth]
\begin{column}{0.49\textwidth}\scriptsize
\textbf{\textcolor{cafe}{A favor de que sea original}}
\begin{itemize}
\item Lectura más difícil: extraño que la misericordia de Israel sea «ahora» tras 11:26.
\item Explica la omisión y el cambio a \grc{ὕστερον}.
\item NET Bible: lo prefiere ligeramente.
\end{itemize}
\textbf{\textcolor{cafe}{A favor de que sea añadido}}
\begin{itemize}
\item P46 y la mayoría no lo tienen.
\item Pudo agregarse para igualar los dos \grc{νῦν} previos.
\item El paralelo pide contraste: «en otro tiempo» / «ahora».
\end{itemize}
\end{column}
\begin{column}{0.49\textwidth}\scriptsize
\begin{tcolorbox}[colback=piedra,colframe=ladrillo,boxrule=0.8pt,arc=2pt,fontupper=\scriptsize]
\textbf{No puede ser base segura de una doctrina.}
\end{tcolorbox}
\textbf{\textcolor{cafe}{Si es original}}
\begin{itemize}
\item «Ahora» escatológico: Ro 3:21; 8:1; 13:11; 2~Co 6:2.
\item En 11:5, el «ahora» es el tiempo \textbf{del remanente}.
\item \grc{ἵνα} + subjuntivo aoristo: \textbf{propósito} de Dios, no cumplimiento total.
\item Lo que dura «ahora» es la desobediencia; el endurecimiento dura «\textbf{hasta que}» (11:25).
\end{itemize}
\end{column}
\end{columns}
\end{frame}

\begin{frame}{Dos predicaciones sobre el «ahora»}
\begin{columns}[T,onlytextwidth]
\begin{column}{0.49\textwidth}\scriptsize\RaggedRight
\textbf{\textcolor{ladrillo}{Michelén (según el resumen)}}\par\smallskip
El «ahora» de 11:31 indica un proceso presente y continuo: «todo Israel» = suma de los elegidos.\par\smallskip
\textbf{Postura 2.}
\end{column}
\begin{column}{0.49\textwidth}\scriptsize\RaggedRight
\textbf{\textcolor{ladrillo}{«El misterio revelado» (Ro 11:25--32)}}\par\smallskip
El endurecimiento «no es total ni es final»; al entrar la plenitud de los gentiles habrá una conversión masiva de Israel. El «ahora» de la misericordia para Israel mira a la \textbf{etapa final}.\par\smallskip
\textbf{Postura 6.}
\end{column}
\end{columns}\par\medskip
\begin{tcolorbox}[colback=piedra,colframe=ladrillo,boxrule=1pt,arc=3pt,fontupper=\small,width=0.95\linewidth]
Las dos coinciden en el \textbf{vaivén} de la misericordia. Difieren en si el «ahora» cierra el futuro o lo anticipa; el texto (\grc{ἄχρι οὗ}, \grc{σωθήσεται}, \grc{ἥξει}) favorece la segunda.
\end{tcolorbox}
\end{frame}

\begin{frame}{Conclusión: ya y todavía no}
\begin{columns}[T,onlytextwidth]
\begin{column}{0.5\textwidth}\scriptsize
\centering\renewcommand{\arraystretch}{1.12}
\begin{tabular}{>{\RaggedRight\arraybackslash}p{1.35cm}>{\RaggedRight\arraybackslash}p{2.05cm}>{\RaggedRight\arraybackslash}p{2.05cm}}
\toprule
\textbf{Pasaje} & \textbf{«Ya»} & \textbf{«Todavía no»}\\\midrule
11:1--6 & Remanente & ---\\
11:7--10 & Endurecidos & Su límite\\
11:11--12 & Celos & Plenitud\\
11:13--16 & «Algunos» & Aceptación\\
11:17--24 & Injerto gentil & Reinjerto\\
11:25--27 & Endurecimiento parcial & Todo Israel; el Libertador\\
11:28--32 & Misericordia a naciones & Misericordia a Israel\\
\bottomrule
\end{tabular}
\end{column}
\begin{column}{0.48\textwidth}\small
\begin{enumerate}
\item \textbf{Sí hay vaivén presente}: la misericordia a las naciones es el medio de la misericordia a Israel.
\item \textbf{El «ahora» no elimina el futuro}: es incierto en los manuscritos y, si es original, es el «ahora» mesiánico que incluye el remanente.
\end{enumerate}
\begin{tcolorbox}[colback=piedra,colframe=ladrillo,boxrule=1pt,arc=3pt,fontupper=\scriptsize]
Michelén capta la columna del «ya»; la postura 6 integra las dos.
\end{tcolorbox}
\end{column}
\end{columns}
\end{frame}

% ===================================================
\section{«Todo Israel»}
\begin{frame}{Seis lecturas de «todo Israel»}
\centering
\begin{tikzpicture}[level distance=1.25cm,
  level 1/.style={sibling distance=6.6cm},
  level 2/.style={sibling distance=1.9cm},
  every node/.style={caja,align=center,text width=1.6cm,font=\tiny}]
\node[text width=2.6cm,font=\scriptsize\bfseries,fill=oro!25] {«Todo Israel» (Ro 11:26)}
  child {node[text width=2.6cm,font=\scriptsize\bfseries] {No étnico}
    child {node {1. La iglesia (Israel espiritual)}}
    child {node {5. Dos Casas (Judá + «Efraín» gentil)}}}
  child {node[text width=2.6cm,font=\scriptsize\bfseries] {Israel étnico}
    child {node {2. Remanente de todas las épocas}}
    child {node {3. Todos los judíos}}
    child {node {4. Dos caminos (sin Cristo)}}
    child {node[fill=oro!25] {6. Conjunto futuro que reconoce al Mesías}}};
\end{tikzpicture}\par\vspace{6pt}
{\scriptsize Organización basada en el material de Danilo Gómez Correa (s.\,f.), con la postura de los «dos caminos» añadida.}
\end{frame}

\begin{frame}{1. La iglesia como Israel}
\begin{columns}[T,onlytextwidth]
\begin{column}{0.49\textwidth}\small\RaggedRight
\textbf{\textcolor{cafe}{Exposición}}\par
«Extiendo el nombre de Israel a todo el pueblo de Dios» (Calvino, 1540/1849). Wright (2013): Israel redefinido en torno al Mesías. Coalición por el Evangelio (en la Clase 7): las promesas se cumplen en la iglesia. Reemplazo estricto: Justino (\emph{Diálogo} 11, 135).\par\smallskip
\textbf{\textcolor{cafe}{Fortaleza}}: unidad en Cristo (Ef 2:14--16; Gá 3:28--29).
\end{column}
\begin{column}{0.49\textwidth}\scriptsize\RaggedRight
\textbf{\textcolor{ladrillo}{Debilidades}}
\begin{itemize}
\item En 11:25 «Israel» se distingue de «las naciones»: cambiar el sentido en 11:26 es improbable.
\item El «cuánto más» (11:12, 15) pierde sentido.
\item La irrevocabilidad se dice del Israel \emph{no creyente} (11:28--29).
\item Desarraiga Jer 31:35--37; Ez 36:24--28; Am 9:14--15.
\item El olivo enseña injerto, no reemplazo (11:17, 24).
\end{itemize}
\end{column}
\end{columns}
\end{frame}

\begin{frame}{2--4. Remanente, todos, dos caminos}
\begin{columns}[T,onlytextwidth]
\begin{column}{0.32\textwidth}\scriptsize\RaggedRight
\textbf{\textcolor{ladrillo}{2. Remanente de todas las épocas}}\par\smallskip
Hendriksen (1981); Robertson (2000); Michelén (s.\,f.).\par\smallskip
\textbf{A favor}: «Israel» étnico; doctrina del remanente.\par
\textbf{En contra}: 11:26 repetiría 11:5; no habría «misterio»; ignora «hasta que» y la «plenitud» frente al remanente.
\end{column}
\begin{column}{0.32\textwidth}\scriptsize\RaggedRight
\textbf{\textcolor{ladrillo}{3. Todos los judíos}}\par\smallskip
Sin respaldo exegético serio.\par\smallskip
\textbf{En contra}: Ro 9:6--8; 10:1 (Pablo ora por su salvación); 11:23 («si no permanecen en incredulidad»). La Mishná misma pone excepciones.
\end{column}
\begin{column}{0.33\textwidth}\scriptsize\RaggedRight
\textbf{\textcolor{ladrillo}{4. Dos caminos (\emph{Sonderweg})}}\par\smallskip
Stendahl (1976); Gaston (1987); Gager (2000): Israel salvado por su pacto, sin fe en Jesús.\par\smallskip
\textbf{A favor}: rechaza el desprecio cristiano a Israel.\par
\textbf{En contra}: Ro 10:9--13 («no hay diferencia»); 11:23; el Libertador es Jesús (1~Ts 1:10); Fil 3:4--9.
\end{column}
\end{columns}
\end{frame}

\begin{frame}{5. Las Dos Casas}
\begin{columns}[T,onlytextwidth]
\begin{column}{0.47\textwidth}\small\RaggedRight
\textbf{\textcolor{cafe}{Exposición}} (Koniuchowsky, 2003)\par
Efraín, dispersado por Asiria (722 a.\,e.\,c.), estaría entre las naciones; la «plenitud de los gentiles» sería la \heb{מְלֹא־הַגּוֹיִם} de Gn 48:19; «todo Israel» = Judá + Efraín (Ez 37:15--28; Os 1:10--11 [2:1--2]).\par\smallskip
\textbf{\textcolor{cafe}{Fortalezas}}: rechaza el reemplazo; toma en serio Ez 37 y Zac 12:10.
\end{column}
\begin{column}{0.51\textwidth}\scriptsize\RaggedRight
\textbf{\textcolor{ladrillo}{Debilidades}}
\begin{enumerate}
\item Los gentiles son \grc{κατὰ φύσιν ἀγριέλαιος}, silvestres \emph{por naturaleza} (11:24).
\item Inclusión por fe, no por linaje (Ro 9:24; Ef 2:11--13).
\item Os 2:23 en Ro 9:25--26: analogía, no genealogía.
\item LXX de Gn 48:19: \grc{πλῆθος}, no \grc{πλήρωμα}.
\item Muchos del norte se integraron en Judá (2~Cr 11:13--17; 15:9; 30:1--11); Ana, de Aser (Lc 2:36); «nuestras doce tribus» (Hch 26:7).
\item Sin evidencia que identifique pueblos gentiles con las tribus.
\end{enumerate}
\end{column}
\end{columns}
\end{frame}

\begin{frame}{6. Israel étnico futuro}
\begin{columns}[T,onlytextwidth]
\begin{column}{0.49\textwidth}\small\RaggedRight
\textbf{\textcolor{cafe}{Exposición}}\par
El pueblo judío como conjunto, no cada individuo, reconocerá a Jesús como Mesías cuando entre la plenitud de las naciones: Murray (1965); Cranfield (1979); Dunn (1988); Morris (1988); Fitzmyer (1993); Stott (1994); Moo (1996); Schreiner (1998); ya Moule (1894).\par\smallskip
{\scriptsize Murray (reformado): es \textbf{exegéticamente imposible} dar a «Israel» en 11:26 otro sentido que el del capítulo; Stott lo suscribe.}\par\smallskip
\textbf{\textcolor{cafe}{Debilidad}}: depende de un cumplimiento futuro no detallado; relación colectivo-individual (el colectivo está hecho de individuos que creen).
\end{column}
\begin{column}{0.49\textwidth}\scriptsize\RaggedRight
\textbf{\textcolor{ladrillo}{Fortalezas}}
\begin{itemize}
\item Mismo sentido de «Israel» en 11:25 y 11:26.
\item Explica «hasta que», remanente/plenitud y «cuánto más».
\item Reinjerto condicionado a la fe (11:23).
\item Uso colectivo de \emph{kol Yisra'el}.
\item Is 59 y Jer 31: pacto con Jacob/Israel.
\item Jesús: «no me veréis más hasta que digáis: Bendito el que viene» (Mt 23:39); Lc 21:24.
\end{itemize}
\end{column}
\end{columns}
\end{frame}

\begin{frame}{Tabla comparativa}
\centering\tiny\renewcommand{\arraystretch}{1.25}
\begin{tabular}{>{\RaggedRight\arraybackslash}p{3.4cm}>{\centering\arraybackslash}p{0.95cm}>{\centering\arraybackslash}p{0.95cm}>{\centering\arraybackslash}p{0.95cm}>{\centering\arraybackslash}p{0.95cm}>{\centering\arraybackslash}p{0.95cm}>{\centering\arraybackslash}p{0.95cm}}
\toprule
\textbf{Criterio (Ro 11)} & \textbf{1. Iglesia} & \textbf{2. Rema\-nente} & \textbf{3. Todos} & \textbf{4. Dos caminos} & \textbf{5. Dos Casas} & \textbf{6. Israel futuro}\\\midrule
Mismo «Israel» en 11:25 y 26 & no & sí & sí & sí & no & \textbf{sí}\\
Explica «hasta que» & no & no & parcial & sí & sí & \textbf{sí}\\
Plenitud mayor que remanente (11:12, 15) & no & no & sí & sí & parcial & \textbf{sí}\\
Exige fe (10:9--13; 11:23) & sí & sí & no & no & sí & \textbf{sí}\\
Gentiles, olivo silvestre (11:24) & parcial & sí & sí & sí & no & \textbf{sí}\\
Irrevocable para el Israel no creyente (11:28--29) & no & parcial & sí & sí & parcial & \textbf{sí}\\
\bottomrule
\end{tabular}\par\vspace{6pt}
\scriptsize\renewcommand{\arraystretch}{1.2}
\begin{tabular}{>{\RaggedRight\arraybackslash}p{2.4cm}>{\RaggedRight\arraybackslash}p{4.0cm}>{\RaggedRight\arraybackslash}p{4.4cm}}
\toprule
\textbf{Postura} & \textbf{Fortaleza} & \textbf{Debilidad}\\\midrule
Israel espiritual & Unidad en Cristo & Ignora la dimensión literal y nacional\\
Dos Casas & Cumplimiento literal; Judá e Israel & Gentiles como Efraín: especulativo\\
Israel étnico con Mesías & Literal y espiritual; integra a los gentiles & Depende de un cumplimiento futuro\\
\bottomrule
\end{tabular}\par
{\fontsize{5}{6}\selectfont\color{sepia}\raggedleft\textit{Nota.} Criterios de elaboración propia; segunda tabla adaptada de Gómez Correa (s.\,f.).\par}
\end{frame}

\begin{frame}{Conclusión sobre «todo Israel»}
\begin{tcolorbox}[colback=piedra,colframe=ladrillo,boxrule=1pt,arc=3pt,fontupper=\small]
La lectura más consistente con el contexto de Romanos 11 es la \textbf{restauración futura del Israel étnico como conjunto, que reconocerá al Mesías}.
\end{tcolorbox}
\begin{columns}[T,onlytextwidth]
\begin{column}{0.44\textwidth}\small
\begin{itemize}
\item \textbf{No cada individuo}: colectivo, como en el AT y la Mishná.
\item \textbf{No por etnia sin fe}: «si no permanecen en incredulidad» (11:23).
\item \textbf{No en reemplazo de la iglesia}: un solo olivo.
\end{itemize}
\end{column}
\begin{column}{0.53\textwidth}
\begin{cita}[fontupper=\scriptsize\itshape]«Todo Israel» significa un futuro volverse masivo del pueblo judío a Yeshúa como Mesías, no la salvación de cada judío individual. Israel es ampliado, no reemplazado.\fuente{(ONE FOR ISRAEL, 2025; síntesis)}\end{cita}
\scriptsize Estimación del ministerio: c. 1 millón de judíos creyentes de c. 15 millones: un remanente.
\end{column}
\end{columns}
\end{frame}

\begin{frame}{Lo que queda abierto}
\begin{columns}[T,onlytextwidth]
\begin{column}{0.32\textwidth}\small\RaggedRight
\textbf{\textcolor{ladrillo}{El momento}}\par\smallskip
Pablo no fecha la plenitud ni la salvación de Israel. «No os toca saber los tiempos» (Hch 1:7).
\end{column}
\begin{column}{0.33\textwidth}\small\RaggedRight
\textbf{\textcolor{ladrillo}{La parusía}}\par\smallskip
¿El Libertador de Sion es la segunda venida? Zac 12:10 («mirarán a mí, a quien traspasaron»); Mt 23:39; Hch 3:19--21 (\grc{ἀποκατάστασις πάντων}).
\end{column}
\begin{column}{0.32\textwidth}\small\RaggedRight
\textbf{\textcolor{ladrillo}{La tierra}}\par\smallskip
Ro 11 no la menciona; los profetas sí (Ez 36--37; Jer 31:38--40). 1948 como Ez 37:1--8 es interpretación teológica de la congregación, no texto de Romanos.
\end{column}
\end{columns}\vspace{8pt}
\centering
\begin{tikzpicture}[node distance=0.25cm]
\foreach \i/\t/\r in {1/{Se juntan los huesos}/{Ez 37:1--7},2/{Cuerpo sin espíritu}/{Ez 37:8},3/{Espíritu: «mirarán»}/{Zac 12:10--14},4/{Un solo pueblo}/{Ez 37:15--25}}{
  \node[caja,text width=2.3cm,align=center] (e\i) at (\i*2.85,0) {\textbf{\textcolor{ladrillo}{\i. \t}}\\{\tiny\color{sepia}\r}};}
\foreach \a/\b in {1/2,2/3,3/4}{\draw[-{Latex[length=2mm]},oro!90!ladrillo,line width=1.2pt] (e\a) -- (e\b);}
\end{tikzpicture}\par
{\fontsize{5}{6}\selectfont\color{sepia}\raggedleft\textit{Nota.} Las cuatro etapas de Ez 37, según la Clase 7 y la Octava lección de la congregación.\par}
\end{frame}

% ===================================================
\section{Hilos}
\begin{frame}{Siete estudios, un capítulo}
\centering\tiny\renewcommand{\arraystretch}{1.2}
\begin{tabular}{>{\RaggedRight\arraybackslash}p{2.0cm}>{\RaggedRight\arraybackslash}p{3.4cm}>{\RaggedRight\arraybackslash}p{2.4cm}>{\RaggedRight\arraybackslash}p{3.4cm}}
\toprule
\textbf{Estudio} & \textbf{Hilo} & \textbf{Textos} & \textbf{En Romanos 11}\\\midrule
1. Abraham--Éxodo & Pacto perpetuo; remanente de José & Gn 15; 17:7; 45:7 & Raíz santa (11:16); amados por los padres (11:28)\\
2. Éxodo--Josué & Reino de sacerdotes; celos & Éx 19:6; Dt 29:4; 32:21 & Celos (11:11, 14); ojos que no ven (11:8)\\
3. Jueces--David & YHWH no desampara; Benjamín & 1~S 12:22 & «¿Ha desechado Dios?» (11:1)\\
4. David--división & Pacto davídico; raíz de Isaí & 2~S 7; Is 11:10 & David testigo (11:9); Ro 15:12\\
5. Reyes y cautiverios & Elías en el norte; «no pueblo mío»; nuevo pacto; Ez 37 & 1~R 19:18; Os 1--2; Jer 31 & Siete mil (11:4); pacto (11:27); irrevocable (11:29)\\
6. Segundo Templo & Redentor de Sion; «todo Israel» confiesa & Is 59:20; Dn 9:11 & «Vendrá de Sion» (11:26)\\
7. Cristo--año 70 & Remanente; «hasta que» & Hch 21:20; Lc 21:24; Mt 23:39 & Remanente (11:5); «hasta que» (11:25)\\
\bottomrule
\end{tabular}
\end{frame}

\begin{frame}{Tres hilos mayores}
\begin{columns}[c,onlytextwidth]
\begin{column}{0.50\textwidth}
\centering
\begin{tikzpicture}[scale=0.82,transform shape]
\node[caja,fill=oro!25,text width=2.4cm,align=center,font=\scriptsize\bfseries] (r) at (0,0) {\textcolor{ladrillo}{Romanos 11}};
\foreach \ang/\n/\t in {90/1/{1: padres y pacto},38.6/2/{2: sacerdocio y celos},-12.9/3/{3: no desamparará},-64.3/4/{4: David y la raíz},-115.7/5/{5: remanente, nuevo pacto},-167.1/6/{6: Redentor de Sion},141.4/7/{7: «hasta que»}}{
  \node[caja,text width=1.9cm,align=center,font=\tiny] (n\n) at (\ang:2.9) {Est. \t};
  \draw[-{Latex[length=2mm]},oro!90!ladrillo,line width=1pt] (n\n) -- (r);}
\end{tikzpicture}
\end{column}
\begin{column}{0.47\textwidth}\small
\begin{enumerate}
\item \destacar{Remanente}: José (Gn 45:7), Elías (1~R 19:18), «un resto volverá» (Is 10:21), el retorno (Esd 9:8), Hch 21:20, Pablo.
\item \destacar{Reunificación}: Judá y José «un solo palo» (Ez 37:15--28). En tiempos de Pablo, «Israel» ya es el heredero de todas las tribus.
\item \destacar{Nuevo pacto}: Jer 31; Ez 36 $\to$ Lc 22:20 $\to$ Ro 11:27.
\end{enumerate}
\end{column}
\end{columns}
\end{frame}

% ===================================================
\section{Tradición judía}
\begin{frame}{La Mishná: «todo Israel»}
\begin{columns}[T,onlytextwidth]
\begin{column}{0.52\textwidth}
\centering{\heb{כל ישראל יש להם חלק לעולם הבא}}\par\smallskip
\begin{cita}[fontupper=\scriptsize\itshape]Todo Israel tiene parte en el mundo venidero, porque está dicho: «Tu pueblo, todos ellos justos» (Is 60:21). Y estos no tienen parte: el que niega la resurrección [\ldots], el que dice que la Torá no es del cielo, y el epicúreo.\fuente{(Mishná, Sanedrín 10:1)}\end{cita}
\scriptsize 10:2 añade Jeroboam, Acab, Manasés, Balaam, Doeg, Ajitófel y Guejazí. \destacar{«Todo Israel» es colectivo con excepciones.}
\end{column}
\begin{column}{0.45\textwidth}\small\RaggedRight
\textbf{\textcolor{cafe}{Las diez tribus (Sanedrín 10:3)}}\par\smallskip
R. Aquiba: «las diez tribus no volverán» (Dt 29:28 [29:27]).\par\smallskip
R. Eliézer: como el día que se oscurece y vuelve a aclarar, así vendrá la luz para ellas.\par\medskip
\precision{Mishná c. 200 e.\,c.: posterior a Pablo. Algunos manuscritos no traen la frase inicial; verificar en edición crítica. Is 60:21 sigue a Is 59:20--21.}
\end{column}
\end{columns}
\end{frame}

\begin{frame}{Talmud, Targum y liturgia}
\begin{columns}[T,onlytextwidth]
\begin{column}{0.49\textwidth}\scriptsize\RaggedRight
\textbf{\textcolor{ladrillo}{Sanedrín 97b--98a}} (tradiciones posteriores)
\begin{itemize}
\item R. Eliézer: «si Israel se arrepiente, es redimido».
\item Is 60:22: «si lo merecen, lo apresuraré; si no, a su tiempo» (R. Josué ben Leví).
\item El Mesías responde «hoy»: Sal 95:7, «si oís su voz».
\item R. Yojanán: angustias como un río $\to$ «vendrá a Sion el Redentor» (Is 59:19--20).
\end{itemize}
\textbf{\textcolor{ladrillo}{Targum de Is 59:20}}: el redentor viene para hacer volver a la Torá a los rebeldes de Jacob (Chilton, 1987).
\end{column}
\begin{column}{0.49\textwidth}\scriptsize\RaggedRight
\textbf{\textcolor{ladrillo}{\emph{U-va le-Tsiyyón}}}: Is 59:20--21 abre una oración del servicio diario de la mañana.\par\smallskip
\textbf{\textcolor{ladrillo}{La Amidá}}\par\smallskip
\renewcommand{\arraystretch}{1.2}
\begin{tabular}{>{\RaggedRight\arraybackslash}p{1.3cm}>{\RaggedRight\arraybackslash}p{3.6cm}}
\toprule
10.ª & «Reúnes a los dispersos de tu pueblo Israel» (Is 11:12; Ez 37:21)\\
14.ª & «Edificas Jerusalén» (Lc 21:24)\\
15.ª & «Haz brotar el renuevo de David» (Jer 23:5; Ro 15:12)\\
\bottomrule
\end{tabular}
\end{column}
\end{columns}
\end{frame}

\begin{frame}{Israel como olivo}
\begin{columns}[c,onlytextwidth]
\begin{column}{0.52\textwidth}
\centering{\large\heb{זַיִת רַעֲנָן יְפֵה פְרִי־תֹאַר}}\par\smallskip
\begin{cita}Olivo verde, hermoso en fruto y en aspecto, llamó YHWH tu nombre; [\ldots] y quebraron sus ramas.\fuente{(Jer 11:16)}\end{cita}
\end{column}
\begin{column}{0.45\textwidth}
\begin{cita}[fontupper=\scriptsize\itshape]¿Por qué Israel es comparado con el olivo? Así como las hojas del olivo no caen ni en verano ni en invierno, así Israel no será desechado nunca, ni en este mundo ni en el venidero.\fuente{(Talmud, Menajot 53b, R. Josué ben Leví)}\end{cita}
\end{column}
\end{columns}\vspace{6pt}
{\small\itshape\color{sepia}\centering El olivo puede perder ramas, pero no es arrancado: la imagen de Pablo habla el idioma de Israel.\par}
\end{frame}

% ===================================================
\section{Implicaciones}
\begin{frame}{Para la iglesia de las naciones}
\centering
\begin{tikzpicture}
 \foreach \x/\n/\t/\d/\r in {0/1/{Humildad}/{«No te jactes contra las ramas». La historia del desprecio cristiano a Israel debe reconocerse como pecado.}/{Ro 11:18},
   3.15/2/{Gratitud}/{Somos «deudores»: Escritura, pactos, Mesías y apóstoles vinieron de Israel.}/{Ro 15:27; 9:4--5},
   6.3/3/{Provocar celos}/{Que Israel vea en la iglesia la bendición que es suya.}/{Ro 11:11, 14},
   9.45/4/{Temor}/{Estamos en el olivo por la fe; la advertencia sigue vigente.}/{Ro 11:20--22}}{
  \node[caja,anchor=north] (b\n) at (\x,0) {\parbox[t][2.9cm][t]{2.65cm}{\RaggedRight
    {\small\textbf{\textcolor{ladrillo}{\t}}}\\[4pt]\scriptsize \d\par\vfill{\tiny\color{sepia}(\r)}}};
  \node[circle,fill=cafe,text=piedra,font=\scriptsize\bfseries,inner sep=1.6pt] at (b\n.north west) {\n};}
\end{tikzpicture}
\end{frame}

\begin{frame}{Contra el desprecio a Israel}
\begin{columns}[c,onlytextwidth]
\begin{column}{0.55\textwidth}
\begin{cita}[fontupper=\small\itshape]Qué extraño que Dios eligiera a los judíos.\fuente{(atribuido a W. N. Ewer)}\end{cita}
\begin{cita}[fontupper=\small\itshape]Pero no tan extraño como los que eligen al Dios judío y desprecian a los judíos.\fuente{(atribuido a C. Browne)}\end{cita}
\end{column}
\begin{column}{0.42\textwidth}\scriptsize\RaggedRight
\begin{itemize}
\item «No te jactes contra las ramas» (Ro 11:18).
\item «De los cuales, según la carne, vino Cristo» (Ro 9:5): el Libertador viene de Sion.
\item «La salvación viene de los judíos» (Jn 4:22).
\item Jonás se enojó porque Dios perdonó a Nínive (Jon 4:1--2): el gentil que se enoja por la misericordia a Israel repite su error.
\end{itemize}
\end{column}
\end{columns}
\end{frame}

\begin{frame}{Israel hoy y la esperanza}
\begin{columns}[T,onlytextwidth]
\begin{column}{0.49\textwidth}\small\RaggedRight
\textbf{\textcolor{cafe}{Teología, no política}}\par\smallskip
Ro 11 habla del pueblo de Israel ante Dios y el Mesías; no es un programa político. Amar a Israel «por causa de los padres» no obliga a aprobar cada acción de un gobierno: los profetas amaban a Israel y denunciaban a sus reyes. Pero Ro 11 impide la indiferencia (Jer 31:35--37).
\end{column}
\begin{column}{0.49\textwidth}
\textbf{\small\textcolor{cafe}{La esperanza}}\par\smallskip
\begin{cita}[fontupper=\scriptsize\itshape]El anhelo de mi corazón y mi oración a Dios por Israel es para salvación.\fuente{(Ro 10:1)}\end{cita}
\small La historia termina con la plenitud de ambos: de las naciones y de Israel.
\end{column}
\end{columns}
\end{frame}

% ===================================================
\section{Síntesis}
\begin{frame}{Síntesis}
\begin{tcolorbox}[colback=piedra,colframe=ladrillo,boxrule=1pt,arc=3pt,fontupper=\small]
\textbf{\textcolor{cafe}{Romanos 11:}} Dios no ha desechado a su pueblo. Hoy hay un remanente escogido por gracia; el resto está endurecido en parte y por un tiempo, para que entren las naciones y para provocar celos a Israel. Cuando haya entrado la plenitud de las naciones, todo Israel será salvo por el Libertador de Sion y el nuevo pacto, porque los dones y el llamamiento de Dios son irrevocables.
\end{tcolorbox}
\scriptsize\renewcommand{\arraystretch}{1.2}
\begin{tabular}{>{\RaggedRight\arraybackslash}p{2.4cm}>{\RaggedRight\arraybackslash}p{8.9cm}}
\toprule
\textbf{Dimensión} & \textbf{Aporte}\\\midrule
Histórica & Roma tras el edicto de Claudio: Pablo corrige la arrogancia gentil (11:13--25)\\
Literaria & Ro 9--11 es una unidad; cuatro pasos: remanente, tropiezo, misterio, doxología\\
Filológica & \grc{πώρωσις ἀπὸ μέρους}, \grc{ἄχρι οὗ}, \grc{πλήρωμα}, \grc{καὶ οὕτως}, \grc{πᾶς Ἰσραήλ}, \grc{ἐκ Σιών}\\
Teológica & Fidelidad de Dios; inclusión sin reemplazo; salvación futura del Israel étnico por la fe\\
\bottomrule
\end{tabular}
\end{frame}

\begin{frame}{La doxología de la serie}
\fotolado{panel_ro33.jpg}{Elaboración propia: Ro 11:33--34a en escritura uncial (texto NA28).}
\begin{columns}[c,onlytextwidth]
\begin{column}{0.57\textwidth}
\small De Gn 15:5, «mira los cielos y cuenta las estrellas», a un apóstol de la tribu de Benjamín que contempla el plan de Dios:
\begin{cita}¡Oh profundidad de las riquezas, y de la sabiduría y del conocimiento de Dios! ¡Cuán insondables son sus juicios e inescrutables sus caminos! [\ldots] Porque de él, por él y para él son todas las cosas. A él sea la gloria por los siglos. Amén.\fuente{(Ro 11:33, 36)}\end{cita}
\end{column}
\begin{column}{0.41\textwidth}\end{column}
\end{columns}
\end{frame}

\begin{frame}{Para la reflexión}
\small
\begin{enumerate}
\item Pablo prueba la fidelidad de Dios con su propia persona (11:1). ¿Qué significa que la prueba sea alguien del remanente?
\item El endurecimiento es parcial, temporal y con propósito. ¿Cómo cambia esto la manera de hablar del pueblo judío?
\item ¿Qué formas de «jactancia contra las ramas» reconoce en la historia de la iglesia y en su comunidad (11:18--22)?
\item ¿Qué versículos de Ro 11 son decisivos para elegir entre «la iglesia» y «el Israel étnico futuro»?
\item La Mishná dice «todo Israel» y enumera excepciones. ¿Qué enseña sobre la frase de Pablo?
\item ¿Qué significa, en la práctica, ser «deudores» de Israel (Ro 15:27)?
\end{enumerate}
\end{frame}

% ===================================================
\section{Glosario y fuentes}
\begin{frame}{Glosario}
\begin{columns}[T,onlytextwidth]
\begin{column}{0.49\textwidth}\tiny
\textbf{Agriélaios} (\grc{ἀγριέλαιος}): olivo silvestre (11:17, 24).\par\smallskip
\textbf{Ametamélēta} (\grc{ἀμεταμέλητα}): irrevocables (11:29).\par\smallskip
\textbf{Aparchḗ} (\grc{ἀπαρχή}): primicia; \heb{רֵאשִׁית} (11:16).\par\smallskip
\textbf{Enkentrízō} (\grc{ἐγκεντρίζω}): injertar (11:17--24).\par\smallskip
\textbf{Go'el} (\heb{גֹּאֵל}): redentor (Is 59:20).\par\smallskip
\textbf{Jéleq} (\heb{חֵלֶק}): parte en el mundo venidero (Mishná).\par\smallskip
\textbf{Kaì hoútōs} (\grc{καὶ οὕτως}): «y así» (11:26).\par\smallskip
\textbf{Kol Yisra'el} (\heb{כָּל־יִשְׂרָאֵל}): todo Israel, colectivo.\par\smallskip
\textbf{Leîmma} (\grc{λεῖμμα}): remanente (11:5).\par\smallskip
\textbf{Mḕ génoito} (\grc{μὴ γένοιτο}): «¡de ningún modo!».
\end{column}
\begin{column}{0.49\textwidth}\tiny
\textbf{Mystḗrion} (\grc{μυστήριον}): designio revelado (11:25).\par\smallskip
\textbf{Parazēlóō} (\grc{παραζηλόω}): provocar a celos (11:11).\par\smallskip
\textbf{Pâs Israḗl} (\grc{πᾶς Ἰσραήλ}): todo Israel (11:26).\par\smallskip
\textbf{Plḗrōma} (\grc{πλήρωμα}): plenitud (11:12, 25).\par\smallskip
\textbf{Pṓrōsis} (\grc{πώρωσις}): endurecimiento (11:25).\par\smallskip
\textbf{Próslēmpsis} (\grc{πρόσλημψις}): aceptación (11:15).\par\smallskip
\textbf{She'erit} (\heb{שְׁאֵרִית}): remanente (Gn 45:7).\par\smallskip
\textbf{Sōthḗsetai} (\grc{σωθήσεται}): será salvo (11:26).\par\smallskip
\textbf{U-va le-Tsiyyón} (\heb{ובא לציון}): «y vendrá a Sion».\par\smallskip
\textbf{Zayit} (\heb{זַיִת}): olivo; Israel (Jer 11:16).
\end{column}
\end{columns}
\end{frame}

\begin{frame}[allowframebreaks]{Referencias}
\tiny
\refapa{Baxter, A.~G., \& Ziesler, J.~A. (1985). Paul and arboriculture: Romans 11.17--24. \emph{Journal for the Study of the New Testament, 8}(24), 25--32.}
\refapa{Beale, G.~K., \& Carson, D.~A. (Eds.). (2007). \emph{Commentary on the New Testament use of the Old Testament}. Baker Academic.}
\refapa{\emph{Biblia Hebraica Stuttgartensia} (5.ª ed.). (1997). Deutsche Bibelgesellschaft.}
\refapa{Biblical Studies Press. (2019). \emph{NET Bible} (2.ª ed.), nota textual a Romanos 11:31. \url{https://netbible.org}}
\refapa{Bruno, C.~R. (2008). The Deliverer from Zion: The source(s) and function of Paul's citation in Romans 11:26--27. \emph{Tyndale Bulletin, 59}(1), 119--134.}
\refapa{Calvino, J. (1849). \emph{Commentaries on the epistle of Paul the apostle to the Romans} (J. Owen, Trad.). Calvin Translation Society. (Obra original publicada en 1540)}
\refapa{Chilton, B.~D. (1987). \emph{The Isaiah Targum} (The Aramaic Bible, Vol. 11). Michael Glazier.}
\refapa{Coalición por el Evangelio. (s.\,f.). [Artículo sobre la teología del reemplazo]. Citado en Departamento de Educación (s.\,f.).}
\refapa{Cranfield, C.~E.~B. (1979). \emph{A critical and exegetical commentary on the epistle to the Romans} (Vol. 2). T\&T Clark.}
\refapa{Danby, H. (Trad.). (1933). \emph{The Mishnah}. Oxford University Press.}
\refapa{Das, A.~A. (2007). \emph{Solving the Romans debate}. Fortress Press.}
\refapa{Departamento de Educación, Congregación Iglesia Israelita del Nuevo Pacto. (s.\,f.). \emph{Restauración de Israel} (Clase 7) [Presentación].}
\refapa{Dunn, J.~D.~G. (1988). \emph{Romans 9--16} (Word Biblical Commentary, Vol. 38B). Word Books.}
\refapa{Epstein, I. (Ed.). (1935--1948). \emph{The Babylonian Talmud}. Soncino Press.}
\refapa{Fitzmyer, J.~A. (1993). \emph{Romans} (The Anchor Bible, Vol. 33). Doubleday.}
\refapa{Gager, J.~G. (2000). \emph{Reinventing Paul}. Oxford University Press.}
\refapa{Gaston, L. (1987). \emph{Paul and the Torah}. University of British Columbia Press.}
\refapa{Gómez Correa, D. (s.\,f.). \emph{Interpretación de Romanos 11:25--26} [Presentación].}
\refapa{Hays, R.~B. (1989). \emph{Echoes of Scripture in the letters of Paul}. Yale University Press.}
\refapa{Hendriksen, W. (1981). \emph{Exposition of Paul's epistle to the Romans}. Baker.}
\refapa{\emph{Israel, el pueblo de Dios} (Octava lección). (s.\,f.). [Material de estudio]. Congregación Iglesia Israelita del Nuevo Pacto.}
\refapa{Jewett, R. (2007). \emph{Romans: A commentary} (Hermeneia). Fortress Press.}
\refapa{Kaiser, W.~C. (1998). \emph{The Messiah in the Old Testament}. Zondervan.}
\refapa{Koniuchowsky, M.~Y. (2003). \emph{The Restoration Scriptures: True Name Edition}. Your Arms to Israel.}
\refapa{Michelén, S. (s.\,f.). [Exposición de Romanos 11] [Video]. YouTube.}
\refapa{Moo, D.~J. (1996). \emph{The epistle to the Romans}. Eerdmans.}
\refapa{Morris, L. (1988). \emph{The epistle to the Romans}. Eerdmans.}
\refapa{Moule, H.~C.~G. (1894). \emph{The epistle of St. Paul to the Romans}. Hodder \& Stoughton.}
\refapa{Murray, J. (1965). \emph{The epistle to the Romans} (Vol. 2). Eerdmans.}
\refapa{Nanos, M.~D. (1996). \emph{The mystery of Romans: The Jewish context of Paul's letter}. Fortress Press.}
\refapa{\emph{Novum Testamentum Graece} (28.ª ed.). (2012). Deutsche Bibelgesellschaft.}
\refapa{ONE FOR ISRAEL Ministry. (2025, 3 de julio). \emph{What does «all Israel will be saved» mean in Romans 11?} \url{https://www.oneforisrael.org/what-does-all-israel-will-be-saved-mean-in-romans-11/}}
\refapa{Predicación sobre Romanos 11:25--32: «El misterio revelado». (s.\,f.). [Transcripción].}
\refapa{Rahlfs, A., \& Hanhart, R. (Eds.). (2006). \emph{Septuaginta} (ed. rev.). Deutsche Bibelgesellschaft.}
\refapa{Robertson, O.~P. (2000). \emph{The Israel of God: Yesterday, today, and tomorrow}. P\&R Publishing.}
\refapa{Sacks, J. (Trad.). (2009). \emph{The Koren Siddur}. Koren Publishers.}
\refapa{Sanders, E.~P. (1977). \emph{Paul and Palestinian Judaism}. Fortress Press.}
\refapa{Schreiner, T.~R. (1998). \emph{Romans} (Baker Exegetical Commentary on the New Testament). Baker Academic.}
\refapa{Stanley, C.~D. (1992). \emph{Paul and the language of Scripture}. Cambridge University Press.}
\refapa{Stendahl, K. (1976). \emph{Paul among Jews and Gentiles}. Fortress Press.}
\refapa{Stott, J.~R.~W. (1994). \emph{The message of Romans}. InterVarsity Press.}
\refapa{Stowers, S.~K. (1981). \emph{The diatribe and Paul's letter to the Romans}. Scholars Press.}
\refapa{Suetonio. (1914). \emph{The lives of the Caesars} (J.~C. Rolfe, Trad.). Harvard University Press.}
\refapa{van der Horst, P.~W. (2000). «Only then will all Israel be saved»: A short note on the meaning of \emph{kai houtōs} in Romans 11:26. \emph{Journal of Biblical Literature, 119}(3), 521--525.}
\refapa{Wagner, J.~R. (2002). \emph{Heralds of the good news: Isaiah and Paul «in concert» in the letter to the Romans}. Brill.}
\refapa{Wiefel, W. (1991). The Jewish community in ancient Rome and the origins of Roman Christianity. En K.~P. Donfried (Ed.), \emph{The Romans debate} (ed. rev.). Hendrickson.}
\refapa{Wright, N.~T. (2013). \emph{Paul and the faithfulness of God}. Fortress Press.}
\end{frame}
\end{document}
